\documentclass[11pt]{article}
\usepackage[T1]{fontenc}
\usepackage{lmodern,amsmath,amssymb,amsthm,booktabs,array,microtype,xcolor}
\usepackage[letterpaper,margin=1in]{geometry}
\usepackage[colorlinks=true,linkcolor=blue!45!black,urlcolor=blue!45!black]{hyperref}
\newtheorem{theorem}{Theorem}[section]
\newtheorem{lemma}[theorem]{Lemma}
\newtheorem{proposition}[theorem]{Proposition}
\newtheorem{corollary}[theorem]{Corollary}
\newtheorem{remark}[theorem]{Remark}
\newcommand{\PP}{\mathbb P}
\newcommand{\EE}{\mathbb E}
\newcommand{\comp}{\operatorname{comp}}
\setlength{\emergencystretch}{3em}
\title{The sharp second-order cycle bound for cyclic cuts of size at most five}
\author{Collaborative research working draft}
\date{10 November 2026}
\begin{document}
\maketitle
\begin{abstract}
There is an absolute constant $K$ such that every connected loopless cubic multigraph $G$ of cycle rank $r$ with a cyclic edge cut of size at most five satisfies
\[
 c(G)\le 2^{r-1}+Kr2^{r/2}.
\]
The known Horak--Kim cyclic-square family shows that the order $r2^{r/2}$ is necessary within this class.
The proof combines terminal-path saturation, a path--cycle tradeoff, exact cut-state gluing, and a linear bound for pairs of disjoint cycles.
For five-edge cuts, a protected distance-layer argument gives unconditional path inequalities.
A distinct-representative argument in the Petersen graph then gives a fixed deficit below one half whenever no smaller cyclic cut is present.
The unrestricted sharp-order bound remains open.
\end{abstract}

\section*{Status, inputs, and conventions}
These statements have passed independent mathematical audits within the research team.
The exact cut-state identities and the sharpness formulas also have independent finite checks.
The proofs are not formalized in Lean, and publication priority has not been established.
The lower-bound family is due to Horak and Kim; it is not a new construction.

Two earlier results in the shared research programme are inputs.
First, fixed path density above one half implies additive terminal-distance stability and bounded rank after deleting a suitable shortest path; see the accompanying report \nolinkurl{terminal-path-stability.pdf}.
Second, the asymptotic one-half cycle bound gives $c(H)\le(1/2+o(1))2^{\beta(H)}$ for unrestricted graphs.
The separate report \nolinkurl{polynomial-cycle-error.pdf} strengthens that input.
These two results are not re-proved here.

Rank means $|E|-|V|+\operatorname{comp}$.
Paths and cycles are simple; parallel edges remain distinct, and a parallel pair is a cycle.
A cyclic cut has a cycle on each shore.
A normalized $k$-pole has $k$ distinct degree-two ports and degree three at every other vertex.
The first twelve sections prove the four-cut theorem.
The later sections extend it to five-edge cuts.
Earlier, weaker estimates are retained as steps of the proof.

\clearpage
\tableofcontents
\section{Saturated four-terminal blocks}
\label{sec:second-order-four-poles}

A four-terminal block is a simple connected graph $K$ with four specified vertices
$L_1,L_2,R_1,R_2$ of degree two.
Every other vertex has degree three.
A cross path joins some $L_i$ to some $R_j$.
A through pair consists of two vertex-disjoint cross paths that cover the four terminals.
Write $h=\beta(K)$.
Call $K$ saturated if $P_{L_iR_j}(K)=2^h$ for every $i,j\in\{1,2\}$.

\begin{proposition}\label{prop:saturated-four-pole}
Suppose that $K$ is saturated and has a through pair.
Then $K$ consists of two paths on $h+1$ vertices each, together with a perfect matching between them.
Orient each path from its left terminal to its right terminal.
In these orders, the matching permutation is a concatenation of singleton blocks and adjacent transpositions.
If there are $a$ transposition blocks, then the number of through pairs is exactly $2^a$.
In particular, that number is at most $2^{\lfloor(h+1)/2\rfloor}$.
\end{proposition}

\begin{proof}
Fix cross terminals $s,t$.
The edge sets with odd degrees precisely at $s,t$ form an affine space of size $2^h$.
Since $s,t$ have degree two, every such set consists of one simple $s$--$t$ path and disjoint cycles.
The other vertices have degree at most three.
Saturation therefore implies that every cross path meets every cycle in a vertex.

Choose a through pair $P,Q$.
Let $F$ contain all edges outside $E(P)\cup E(Q)$.
Every vertex of $P\cup Q$ has degree one in $F$.
Every other vertex has degree three in $F$.
A cycle in $F$ would avoid $P$, which is impossible.
Thus $F$ is a forest.

No component of $F$ has two leaves on $P$.
Otherwise, its path between those leaves, together with a subpath of $P$, contains a cycle disjoint from $Q$.
The internal vertices of that forest path lie outside $P\cup Q$, since all vertices on $P\cup Q$ have degree one in $F$.
The same argument applies to $Q$.
Thus every component of $F$ has at most two leaves.
A tree whose vertex degrees belong to $\{1,3\}$ has two more leaves than degree-three vertices.
Each component is consequently one edge, with one endpoint on each path.
Hence $P\cup Q$ spans $K$, and $F$ is a perfect matching between $P$ and $Q$.

Both paths have the same number $k$ of vertices.
Write their vertices as $x_1,\ldots,x_k$ and $y_1,\ldots,y_k$, in left-to-right order.
The matching has edges $x_i y_{\pi(i)}$ for a permutation $\pi$.
Counting edges gives $h=k-1$.

We claim that $\pi$ avoids the patterns $231$, $312$, and $321$.
Suppose that $i<j<k'$ and $\pi(k')<\min\{\pi(i),\pi(j)\}$.
The first two matching edges, their top interval, and their bottom interval form a cycle.
The path from $x_k$ to $x_{k'}$, across its matching edge, and then to $y_1$ avoids that cycle.
It is a cross path, a contradiction.
This excludes both $231$ and $321$.
If instead $\pi(i)>\max\{\pi(j),\pi(k')\}$, use the last two matching edges for the cycle.
The path from $x_1$ through the first matching edge to $y_k$ avoids it.
This excludes $312$.

For completeness, these three forbidden patterns force the claimed block form.
If $\pi(1)\ge3$, its entries with values one and two give a $312$ or $321$ pattern.
If $\pi(1)=1$, remove this singleton and apply induction.
If $\pi(1)=2$, the value one must occur second.
Otherwise, the first entry, second entry, and later value one form a $231$ pattern.
Remove this adjacent transposition and apply induction.

The graph is therefore a chain of single-rung blocks and square blocks.
Between successive blocks there are two edges.
Each through pair uses both edges at every such interface.
A single-rung block has one possible local through pairing.
A square block has two.
The choices are independent, so there are exactly $2^a$ through pairs.
Since each square block uses two positions, $a\le\lfloor k/2\rfloor$.
\end{proof}

Conversely, every graph described in the proposition is saturated.
With $a$ square blocks and $c$ single-rung blocks, its rank is $h=2a+c-1$.
For fixed cross terminals, choose one edge at each of the $a+c-1$ interfaces.
Each square then supplies two local paths, while each single-rung block supplies one.
This gives $2^{a+c-1}2^a=2^h$ cross paths.

The proposition gives a barrier for a restricted construction method.
In a cyclic chain of saturated blocks, the odd-interface sector gives the half-density leading term.
The proposition limits the through-state growth of every block with a through pair to the half exponent.
It does not classify saturated blocks with no through pair.
It also does not bound arbitrary graphs or blocks with unsaturated cross-path counts.

\section{Exact bounds for saturated blocks}
\label{sec:second-order-saturated-blocks}

All multigraphs in this section are finite and loopless.
Parallel edges are allowed and are distinguished when paths are counted.
The four-terminal classification in Proposition~\ref{prop:saturated-four-pole}
also holds for such multigraphs.
Its proof forces a spanning pair of paths and a matching between them.
The resulting permutation ladder is simple.

For positive integers $h$, define
\[
 M(2j)=3\cdot2^j-2,\qquad
 M(2j+1)=4\cdot2^j-2.
\]
Thus $M(1)=2$, and $M(h)\le3\cdot2^{h/2}$.

\begin{theorem}\label{thm:saturated-three-pole}
Let $T$ be connected, with three distinct degree-two terminals $a,b,c$.
Every other vertex has degree three.
Suppose that $\beta(T)=h$ and
$P_{ac}(T)=P_{bc}(T)=2^h$.
Then
\[
 P_{ab}(T)\le M(h).
\]
For every $h\ge1$, a simple graph attains equality.
\end{theorem}

\begin{proof}
Saturation implies that every $a$--$c$ path and every $b$--$c$ path meets every cycle.
Indeed, the corresponding affine parity class has $2^h$ members.
Every member is one terminal path together with vertex-disjoint cycles.

We first reduce bridges by induction on $h$.
Every component obtained by deleting a bridge contains a cycle.
To see this, let the component have $n$ vertices and $p$ original terminals.
Its degree sum is $3n-p-1$, so its rank is $(n-p+1)/2$.
Since $n\ge p$, this integer is positive.
Consequently, a bridge must separate $\{a,b\}$ from $c$.
Otherwise, one saturated terminal pair lies on one side and avoids a cycle on the other.

The endpoint of this bridge on the $\{a,b\}$ side cannot be $a$ or $b$.
Such a terminal path to $c$ would immediately cross the bridge and avoid a cycle on its starting side.
The other endpoint cannot be $c$, by the same argument.
Write the bridge as $zw$, with $a,b,z$ on its left side.
After deleting the bridge, the left graph $T_L$ is a three-terminal graph with terminals $a,b,z$.
The right graph $T_R$ has two degree-two terminals $w,c$ and otherwise degree three.
Write their ranks as $h_L,h_R$.
Both ranks are positive, and $h=h_L+h_R$.
Terminal paths factor across the bridge.
The bounds $P_{az}(T_L)\le2^{h_L}$ and $P_{wc}(T_R)\le2^{h_R}$ force equality in both.
The same argument gives $P_{bz}(T_L)=2^{h_L}$.
Also, $P_{ab}(T)=P_{ab}(T_L)$.
Induction and the monotonicity of $M$ prove the result when a bridge exists.

Now suppose that $T$ has no bridge.
A connected bridgeless graph of maximum degree three has no cut vertex.
Otherwise, each component after deleting a cut vertex requires at least two incident edges at that vertex.
Thus $T$ is two-connected.
When $h=1$, degree counting gives three vertices, so $T$ is a triangle.
Assume $h\ge2$.

The vertex $c$ is not adjacent to $a$ or $b$.
For example, an edge $ac$ gives at most
$1+2^{h-1}<2^h$ simple $a$--$c$ paths.
Here, delete that nonbridge edge to bound all paths other than the edge itself.
The two neighbors $x,y$ of $c$ are distinct.
If both incident edges joined $c$ to one vertex, that vertex's remaining edge would be a bridge.

Set $K=T-c$.
This connected four-terminal graph has terminals $a,b,x,y$ and rank $h-1$.
The equalities
\[
 P_{ax}(K)+P_{ay}(K)=2^h,
 \qquad
 P_{bx}(K)+P_{by}(K)=2^h
\]
force each summand to equal $2^{h-1}$.
Thus $K$ is saturated.
It has a through pair.
Otherwise, the vertex form of Menger's theorem gives a single vertex separating
$\{a,b\}$ from $\{x,y\}$ in $K$.
That vertex would disconnect $T$, contrary to two-connectivity.
The statement remains valid when the separating vertex is one of the four terminals.
At least one left terminal and one right terminal remain after its deletion.

Apply Proposition~\ref{prop:saturated-four-pole}.
Suppose that its ladder has $A$ square blocks and $s$ single-rung blocks.
Then $h=2A+s$, and its through-pair count is $2^A$.
An $a$--$b$ path in $K$ reaches one final block, turns there, and returns along the other strand.
A square block has two possible turns and two through choices.
A single-rung block has one turn and one through choice.
The square turns contribute $2(2^A-1)$ in total.
Each single-rung turn contributes at most $2^A$.
Therefore
\[
 P_{ab}(K)\le (s+2)2^A-2.
\]
Paths in $T$ through $c$ correspond exactly to the through pairs in $K$.
Hence
\[
 P_{ab}(T)=P_{ab}(K)+2^A
 \le(s+3)2^A-2.
\]
Subject to $2A+s=h$, this expression is largest when $A=\lfloor h/2\rfloor$.
Indeed, replacing two single-rung blocks by one square multiplies the positive factor by
$2(s+1)/(s+3)>1$ when $s\ge2$.
The resulting bound is $M(h)$.

For even $h$, take only square blocks and join the two right terminals to a new terminal $c$.
For odd $h\ge3$, add one single-rung block at the right end before adding $c$.
All single-rung turns then have their maximum weight, so equality holds.
For $h=1$, use a triangle.
All these examples are simple.
\end{proof}

\begin{lemma}\label{lem:saturated-two-pole-cycles}
Let $N$ be a connected loopless multigraph with two distinct degree-two terminals $u,v$.
Every other vertex has degree three.
If $\beta(N)=b$ and $P_{uv}(N)=2^b$, then
\[
 c(N)\le S(b):=\sum_{j=1}^b\bigl(M(j)-1\bigr)
 \le10\cdot2^{b/2}.
\]
The same bound holds if extra degree-two vertices subdivide edges.
\end{lemma}

\begin{proof}
Suppress unmarked vertices of degree two.
This preserves rank, terminal-path counts, and cycle counts.
No loop can arise in this normalization.
A pendant cycle would avoid every terminal path and contradict saturation.
The rank-one case is a cycle, with one cycle in total.

Suppose that $b\ge2$, and choose an edge $e=ux$ incident to $u$.
The vertices $u,v$ cannot be adjacent, since a direct edge would give at most $1+2^{b-1}$ terminal paths.
This edge is not a bridge.
If it were, both shores would contain cycles by degree counting.
If $v$ lies opposite $u$, every terminal path immediately crosses the bridge and avoids the cycle on the $u$ shore.
If $v$ lies on the $u$ shore, every terminal path avoids the cycle on the opposite shore.

Subdivide $e$ by a new vertex $a$ to obtain a three-terminal graph $T$ with terminals $u,a,v$.
Sliding the initial endpoint across the new edge gives a bijection between
$u$--$v$ paths and $a$--$v$ paths.
If a path starts with $ua$, remove that edge.
Otherwise, prepend $au$.
The degree-two condition at $u$ makes both operations reversible.
Thus both counts equal $2^b$.
Theorem~\ref{thm:saturated-three-pole} gives
\[
 1+c_e(N)=P_{ua}(T)\le M(b).
\]

Deleting $e$ leaves a connected graph of rank $b-1$.
Exactly half of the affine $u$--$v$ parity class avoids $e$.
Every member of that class is a path, so
$P_{uv}(N-e)=2^{b-1}$.
Delete the now-forced initial leaf segment and suppress unmarked degree-two vertices.
These operations preserve the rank and cycle count, and preserve the terminal-path count after moving the initial terminal.
If the rank is positive, the resulting terminals remain distinct and have degree two.
Otherwise, a forced terminal path would be unique, contrary to $2^{b-1}>1$.
A loop or an off-route cyclic component would contradict saturation.
Thus induction applies to the resulting two-terminal graph.
The identity $c(N)=c(N-e)+c_e(N)$ proves the recurrence.

Direct summation gives
\[
 S(2j)=10\cdot2^j-6j-10,
 \qquad
 S(2j+1)=14\cdot2^j-6j-13.
\]
Both expressions are at most $10\cdot2^{b/2}$.
\end{proof}

\begin{corollary}\label{cor:saturated-four-pole-cap}
For a saturated four-terminal graph $K$ of rank $h$, let $D(K)$ count its through pairs.
Then
\[
 \begin{split}
 c(K)&\le10\cdot2^{h/2},\\
 P_{L_1L_2}(K)+D(K)&\le M(h+1),\\
 P_{R_1R_2}(K)+D(K)&\le M(h+1),\\
 D(K)&\le2^{\lfloor(h+1)/2\rfloor}.
 \end{split}
\]
\end{corollary}

\begin{proof}
For the first bound, select one cross-terminal pair and suppress the two unused degree-two terminals.
Apply Lemma~\ref{lem:saturated-two-pole-cycles}.
For the second, add a new vertex $c$ joined to $R_1,R_2$.
The resulting three-terminal graph has rank $h+1$.
Each path count from $L_i$ to $c$ is $2^{h+1}$.
Its $L_1$--$L_2$ paths either avoid $c$ or correspond to through pairs in $K$.
Apply Theorem~\ref{thm:saturated-three-pole}.
The third bound is symmetric.
For the last bound, use Proposition~\ref{prop:saturated-four-pole} when $D(K)>0$.
When $D(K)=0$, it is immediate.
\end{proof}

\begin{lemma}\label{lem:saturated-four-pole-double-cap}
For a saturated four-terminal graph $K$ of rank $h$, let $C(K)$ count pairs
of vertex-disjoint cap paths, one from $L_1$ to $L_2$ and one from $R_1$ to $R_2$.
Then
\[
 C(K)\le3(h+1)2^{(h+1)/2}.
\]
\end{lemma}

\begin{proof}
First suppose that $D(K)>0$.
Use its classified ladder, and give a square block weight two and a single-rung block weight one.
Their total weight is $w=h+1$.
A disjoint cap pair turns its left path in a block $i$ and its right path in a later block.
The two paths cannot turn in the same block, since their local turns would intersect.
Fix the first turn block $i$.
The preceding through choices contribute at most $2^{w_{<i}/2}$.
The turn at block $i$ contributes at most $2^{w_i/2}$.
The right cap in the remaining suffix contributes at most $3\cdot2^{w_{>i}/2}$.
For this last bound, a suffix with $A$ squares and $s$ single rungs has at most $(s+2)2^A-2$ caps.
The inequality follows from $(s+2)2^{-s/2}\le3$.
It also covers a one-rung suffix, while an empty suffix contributes zero.
Thus each first-turn block contributes at most $3\cdot2^{w/2}$.
There are at most $w$ blocks.

Now suppose that $D(K)=0$.
Join $R_1,R_2$ to a new vertex $c$, as in Corollary~\ref{cor:saturated-four-pole-cap}.
The resulting saturated three-terminal graph has a bridge.
Otherwise, the bridgeless case in the proof of Theorem~\ref{thm:saturated-three-pole} would give a through pair in $K$.
Every such bridge separates $\{L_1,L_2\}$ from $c$, and neither endpoint is $c$.
It therefore lies in $K$ and separates its left terminal pair from its right terminal pair.

Each component of $K$ after deleting this bridge has positive rank.
Indeed, the same degree count gives rank $(n-p+1)/2\ge1$ for a side containing $p$ original terminals.
Neither bridge endpoint can be an original terminal.
Otherwise, a cross path from that terminal would avoid a cycle on its own side.
The two sides are thus three-terminal graphs, of ranks $h_L,h_R\ge1$, with $h_L+h_R=h$.
The cross-path counts factor across the bridge.
Their saturation forces both path counts to the bridge endpoint to be saturated on each side.
Hence Theorem~\ref{thm:saturated-three-pole} gives
\[
 C(K)=P_{L_1L_2}(K_L)P_{R_1R_2}(K_R)
 \le M(h_L)M(h_R)\le9\cdot2^{h/2}.
\]
Since $h\ge2$, this is at most $3(h+1)2^{(h+1)/2}$.
\end{proof}

\begin{theorem}\label{thm:saturated-ring-barrier}
Take $k\ge2$ vertex-disjoint saturated four-terminal graphs $K_1,\ldots,K_k$.
Join their right terminal pairs to the next left terminal pairs by two matching edges, cyclically.
The matching at each interface can be straight or crossed.
Let $G$ be the resulting graph and let $r=\beta(G)$.
Then
\[
 2^{r-1}\le c(G)\le2^{r-1}+20r2^{r/2}.
\]
Thus this entire saturated-block construction method cannot improve the half exponent in the cycle excess.
\end{theorem}

\begin{proof}
Write $h_i=\beta(K_i)$ and $w_i=h_i+1$.
Each four-terminal graph has $h_i\ge1$, so $w_i\ge2$.
Set $W=\sum_i w_i=r-1$.
Cut parity makes the number of used edges at every interface have the same parity.

For odd parity, choose one edge at each interface.
Each block then supplies $2^{h_i}$ paths between its selected cross terminals.
Every choice gives one simple cycle.
This sector has exactly
$2^k\prod_i2^{h_i}=2^W$ cycles.

An even-parity cycle with no used interface lies inside one block.
Corollary~\ref{cor:saturated-four-pole-cap} bounds their total by
$10\sum_i2^{h_i/2}\le10k2^{W/2}$.
If all interfaces are used, each block supplies a through pair or a pair of disjoint caps.
Two or more cap-pair blocks produce two or more components, so a single cycle uses at most one.
The all-through choices contribute at most $\prod_iD(K_i)\le2^{W/2}$ cycles.
By Lemma~\ref{lem:saturated-four-pole-double-cap}, the choices with one cap-pair block contribute at most
\[
 \sum_i C(K_i)\prod_{j\ne i}D(K_j)
 \le3W2^{W/2}.
\]

Otherwise, the used interfaces form a nonempty proper cyclic interval.
The two end blocks supply cap paths.
Every inner block supplies a through pair.
For a fixed interval whose occupied blocks form $I$, its contribution is at most
\[
 9\cdot2^{\frac12\sum_{i\in I}w_i}.
\]
Here we used $M(h_i+1)\le3\cdot2^{w_i/2}$ at both ends.
If the interval omits $t$ blocks, then their total weight is at least $2t$.
There are $k$ intervals for each $t\in\{0,\ldots,k-2\}$.
Their combined contribution is therefore at most
\[
 9k2^{W/2}\sum_{t=0}^{k-2}2^{-t}
 \le18k2^{W/2}.
\]
The total excess is at most $(28k+3W+1)2^{W/2}$.
Since $2k\le W$, this is at most $(17W+1)2^{W/2}$,
which is smaller than $20r2^{r/2}$.
\end{proof}

\section{A gap below full terminal-path saturation}
\label{sec:second-order-path-gap}

Paths and cycles are simple.
Parallel edges remain distinct, and a pair of parallel edges can form a cycle.
The term \emph{saturated} means that the terminal-path count equals $2^{\beta(G)}$.

\begin{lemma}\label{lem:full-two-pole-classification}
Let $N$ be a connected loopless multigraph with two distinct degree-two terminals $a,c$.
Every other vertex has degree three.
Suppose that $P_{ac}(N)=2^{\beta(N)}$.
Then its bridges join saturated bridgeless two-terminal pieces in series.
Every such piece is either a pair of parallel terminal edges, or a capped permutation ladder.
The ladder permutation is a concatenation of singleton blocks and adjacent transpositions.
\end{lemma}

\begin{proof}
Every terminal path meets every cycle.
This follows by counting its affine parity class, which has $2^{\beta(N)}$ members.
Each member is one terminal path together with vertex-disjoint cycles.

Every shore of a bridge contains a cycle.
If a shore has $n$ vertices and $p$ original terminals, its rank is $(n-p+1)/2\ge1$.
Hence every bridge separates $a$ from $c$.
Neither original terminal meets a bridge.
A path from such a terminal would immediately cross it and avoid a cycle on its starting shore.
All bridges occur on their route, and no cyclic part can lie off that route.
After deleting the bridges, the attachment vertices of each piece are distinct.
Otherwise, every terminal path would avoid a cycle inside that piece.
The pieces are two-terminal graphs with degree-two attachment vertices and all other degrees three.
Their ranks add, and their terminal-path counts multiply.
Saturation forces every piece to be saturated.

Consider one bridgeless piece.
A connected bridgeless graph of maximum degree three has no cut vertex.
Choose two internally vertex-disjoint terminal paths $P,Q$.
Let $F$ contain the unused edges.
The terminals have degree zero in $F$.
Other vertices of $P\cup Q$ have degree one in $F$, while unused vertices have degree three.
The graph $F$ has no cycle, since such a cycle would avoid $P$.
No forest component can have two leaves on the interior of $P$.
Its joining path would form a cycle with a subpath of $P$, disjoint from $Q$.
The same argument applies to $Q$.
Thus every nonempty component of $F$ is one edge joining the two path interiors.

The two path interiors therefore have the same size $k$, and together they contain all nonterminal vertices.
Their unused edges give a perfect matching.
When $k=0$, the piece consists of two parallel terminal edges.
Otherwise, order both interiors from $a$ to $c$ and write the matching as a permutation $\pi$.
A pattern $231$ or $321$ produces a two-rung cycle on its first two positions.
The third rung joins the top right interval to the bottom left interval outside that cycle.
Together with $a,c$, these intervals give a terminal path disjoint from the cycle.
A pattern $312$ gives the symmetric contradiction using its last two positions for the cycle.
Thus $\pi$ avoids these three patterns.
If its first value is at least three, the later values one and two give a $312$ or $321$ pattern.
If its first value is two, value one must occur second, since an intervening value would give $231$.
Delete the initial singleton or transposition and apply induction.
The common endpoints $a,c$ cap the two ends of the ladder.
\end{proof}

\begin{lemma}\label{lem:one-full-three-pole-gap}
Let $T$ be connected, with three distinct degree-two terminals $a,b,c$ and all other degrees three.
If $\beta(T)=h$ and $P_{ac}(T)=2^h$, then
\[
 P_{bc}(T)=2^h
 \quad\hbox{or}\quad
 P_{bc}(T)\le\frac78\,2^h.
\]
The factor $7/8$ is attained by a simple graph.
\end{lemma}

\begin{proof}
Suppress $b$ to an edge of a saturated two-terminal graph $N$.
Its two neighbors are distinct.
Otherwise, its parallel-edge cycle would be pendant and would avoid every $a$--$c$ path.
Apply Lemma~\ref{lem:full-two-pole-classification}.

Suppose first that $b$ subdivides a bridge.
The earlier pieces have total positive rank, while every $b$--$c$ path avoids them.
The normalized count is therefore at most $1/2$.
The same bound holds if $b$ lies in any piece after the first.
In that case, the later pieces multiply both the local path count and its upper bound by the same power of two.
The earlier pieces contribute at least one additional rank unit only to the denominator.
It remains to consider an edge in the first piece.
A piece consisting of two parallel edges gives a triangle after subdivision, and its local count is saturated.

Now use the capped ladder description of the first piece.
Suppose that it has $k$ ladder positions, so its rank is $k+1$.
For a prefix containing $p$ positions and the left cap $a$, let $A_p$ count paths between its two right attachment vertices.
Set $A_0=1$ for the empty prefix.
For $p\ge1$, this prefix has rank $p$, so $A_p\le2^p$.
A suffix containing $q$ positions and the right cap has $2^q$ paths from either left attachment to its terminal.
Any path from $b$ to the right terminal uses exactly one edge of the suffix's two-edge entrance cut.
Its crossing parity is odd, and a simple path cannot use either edge twice.
Thus the suffix contributes one terminal path, and the following local lists are exhaustive.

First suppose that the subdivided edge is an interface edge or an end-cap edge.
With $p$ positions before it and $q$ after it, the path count is
\[
 (1+A_p)2^q,
 \qquad
 \frac{P_{bc}}{2^{k+1}}=\frac{1+A_p}{2^{p+1}}.
\]
The path first goes right, or first uses a prefix cap and returns along the other interface edge.
The displayed ratio is one when $p=0$, and at most $3/4$ when $p\ge1$.

If the subdivided edge is the rung of a single-rung block, its two initial directions give
$2(1+A_p)2^q$ paths.
Here $k=p+1+q$, so the same normalized ratio applies.

Finally, suppose that the subdivided edge lies in a square block.
Delete that edge, and regard a path through the earlier prefix as one virtual cap between the square's left ports.
There are four local paths to the right ports that avoid this cap.
There are three that use it.
For clarity, label the square ports $u,v$ on the left and $x,y$ on the right, and subdivide $ux$ by $b$.
The cap-free paths are
\[
 bx,\quad bu yv x,\quad bu y,\quad bx v y.
\]
The paths using the virtual cap $uv$ are
\[
 bu v x,\quad bu v y,\quad bx v u y.
\]
Each virtual cap has $A_p$ realizations.
Consequently,
\[
 P_{bc}=(4+3A_p)2^q,
 \qquad
 \frac{P_{bc}}{2^{k+1}}=\frac{4+3A_p}{2^{p+3}}.
\]
This ratio is $7/8$ when $p=0$, and at most $5/8$ when $p\ge1$.
All edge types have now been covered.

For equality, take the capped ladder with one square block and subdivide any square edge.
Its old terminal-path count is eight, and the new count is seven.
The graph is simple.
\end{proof}

\begin{theorem}\label{thm:terminal-path-saturation-gap}
Let $G$ be a connected finite graph or multigraph of rank $b$, with distinct terminals $s,t$.
Then
\[
 P_{st}(G)=2^b
 \quad\hbox{or}\quad
 P_{st}(G)\le\frac{15}{16}\,2^b.
\]
The factor $15/16$ is attained by a simple graph of rank four.
\end{theorem}

\begin{proof}
Induct on $b$.
For $b=0$, the graph is a tree and has one terminal path.
Use the block route between $s,t$.
An off-route cyclic block contributes rank but no terminal choices, so it gives $P_{st}(G)\le2^{b-1}$.
Otherwise, the route-block ranks add and their terminal-path counts multiply.
If there are at least two cyclic route blocks, each has smaller rank.
Induction proves the dichotomy by multiplying their normalized counts.
Bridge blocks contribute one and can be suppressed.
Thus it remains to consider a two-connected loopless block with distinct terminal attachments.

At any vertex $v$ of this block, the incident-edge pattern of a uniform affine terminal-parity vector is uniform subject to its required parity.
Indeed, deleting $v$ leaves a connected graph.
Every compatible incident pattern extends to a parity vector there, with the same number of extensions.
At an internal vertex of degree $d$, a simple path uses zero or two incident edges.
If $d\ge4$, this bounds its density by
\[
 \frac{1+\binom d2}{2^{d-1}}\le\frac78.
\]
At a terminal of degree $d\ge3$, only singleton incident patterns are allowed, giving density at most
$d/2^{d-1}\le3/4$.
Both cases prove the required bound immediately.

Suppress unmarked degree-two vertices.
We may therefore assume that both terminals have degree two and every other vertex has degree three.
The rank-one case is a cycle and is saturated.
Assume $b\ge2$.
If the terminals are adjacent, their direct edge and the remaining paths give
$P_{st}(G)\le1+2^{b-1}\le(3/4)2^b$.
Otherwise, let $x,y$ be the neighbors of $s$.
They are distinct.
If both edges joined $s$ to one vertex, that vertex would be a cut vertex unless the graph were the rank-one digon.

The graph $T=G-s$ is connected and has rank $b-1$.
Its three distinct terminals $x,y,t$ have degree two, and every other vertex has degree three.
Split paths according to their first edge:
\[
 P_{st}(G)=P_{xt}(T)+P_{yt}(T).
\]
Apply induction to both terminal pairs in $T$.
If neither count is full, their sum is at most $(15/16)2^b$.
If one is full, Lemma~\ref{lem:one-full-three-pole-gap} makes the other either full or at most $(7/8)2^{b-1}$.
Thus their sum is either $2^b$ or at most
$(1+7/8)2^{b-1}=(15/16)2^b$.

For equality, use the rank-three three-terminal example at the end of Lemma~\ref{lem:one-full-three-pole-gap}.
Add a new terminal $s$ joined to its two left terminals, and keep its right terminal $t$.
The two branch counts are eight and seven.
The resulting simple graph has rank four and fifteen terminal paths.
\end{proof}

In particular, a sequence of four-terminal blocks whose four cross-path densities tend to one is eventually exactly saturated.
Thus vanishing cross-path defects cannot bypass Theorem~\ref{thm:saturated-ring-barrier}.

% Independent audit: path_stability passed the geometric proof and its refinements.
% Exact checks: check_high_path_cycle_layers.py; high-path-cycle-layers-results.json.
\section{Few internal cycles at fixed terminal-path density}
\label{sec:high-path-cycle-tradeoff}

A \emph{normalized two-pole} is a connected loopless multigraph with two distinct degree-two terminals.
Every other vertex has degree three.
Its rank $b$ and order $n$ satisfy $n=2b$.
Parallel edges remain distinct.

\begin{lemma}[Counting cycles through thin distance layers]
\label{lem:thin-layer-cycle-count}
Let $N$ be a bridgeless normalized two-pole of rank $b$, with terminals $s,t$.
Choose a shortest terminal path $A$, of length $\ell$.
Put
\[
 \Delta=b-\ell,\qquad H=\beta(N-V(A)).
\]
Then $\Delta\ge0$, and
\[
 c(N)\le 2^{19\Delta+12H+24+b/2}.
\]
If $N$ is simple, the term $12H$ can be omitted.
\end{lemma}

\begin{proof}
A connected bridgeless graph of maximum degree three has no cut vertex.
Let $L_i$ be its distance-$i$ layer from $s$.
For $1\le i<\ell$, the layer $L_i$ separates $s$ from $t$.
Thus $|L_i|\ge2$.
The endpoint layers each contain at least one vertex.
Consequently, $2b=|V(N)|\ge2\ell$, proving $\Delta\ge0$.

Mark all vertices outside the internal layers $L_1,\ldots,L_{\ell-1}$.
Also mark every internal layer containing at least three vertices.
The total excess above two vertices per internal layer, and one per endpoint layer, is $2\Delta$.
For $w\ge3$, one has $w\le3(w-2)$.
It follows that the marked set has at most $6\Delta+2$ vertices.
Every remaining layer has exactly two vertices, including its vertex of $A$.

Each vertex of an unmarked layer has a neighbor in the preceding layer and in the following layer.
Otherwise, deleting the other vertex of its layer would separate $s$ from $t$.
In particular, each such vertex has at most two incidences towards either neighboring layer.

If $N$ has parallel edges, mark every layer containing an endpoint of a parallel pair outside $V(A)$.
There are at most $H$ such pairs.
Indeed, their digons have disjoint edge supports and are independent in the cycle space of $N-V(A)$.
Parallel pairs have disjoint vertex sets, since the maximum degree is three.
A triple edge would form an isolated cubic component, which is impossible here.
No parallel pair has just one endpoint on $A$, because each vertex of $A$ has at most one incidence outside $A$.
Every parallel pair with both endpoints on $A$ joins consecutive path vertices, by geodesicity.
Each newly marked layer had two vertices, so at most $4H$ vertices are added.
The marked set now has at most $6\Delta+4H+2$ vertices.

Fix a cycle $C$ and its minimum and maximum distance layers $i\le j$.
Mark their vertices as well, unless their layers were already marked.
This adds at most four vertices.
Write $B$ for the final marked set, so
\[
 |B|\le6\Delta+4H+6.
\]
Fix all edges of $C$ incident with $B$.
There are at most $2^{3|B|}$ choices.

Call each unmarked layer \emph{ordinary}.
Consider an ordinary layer strictly between the two extreme layers of $C$.
The cycle must use both vertices of that layer.
Each vertex has one cycle edge to the preceding layer and one to the following layer.
To see this, every prefix cut meets a cycle evenly.
A layer with two vertices therefore carries two strands.
A same-layer edge would leave an odd prefix crossing.
Two separate caps at that layer would disconnect the cycle into a left part and a right part.
Ordinary layers outside the two extreme layers contribute no cycle edges.

Thus, across consecutive ordinary layers used by $C$, its edges form a perfect matching.
Such a cut has at most four edges.
For a simple cut, its perfect matching is unique unless the cut is $K_{2,2}$.
That cut has two perfect matchings.
The only other two-choice cut consists of a doubled edge of $A$ and one edge on the other strand.
Indeed, a doubled cross-strand edge would give an $A$ vertex two incidences outside $A$.
A doubled edge on the other strand has a marked endpoint.
A four-edge cut containing parallel edges must contain two doubled pairs, so it also has a marked endpoint.

Two cuts with two choices cannot be consecutive.
A $K_{2,2}$ cut already uses two incidences at each adjacent vertex.
An adjacent two-choice cut would force degree at least four at an $A$ vertex.
Two consecutive doubled edges of $A$ also force degree four.
Therefore each two-choice cut uses a disjoint pair of two-vertex layers.
There are at most $j-i+1$ such layers, and hence at most $(j-i+1)/2$ two-choice cuts.

The extreme layers, the edges incident with $B$, and these matching choices determine $C$ completely.
Let $D$ be the largest distance layer in $N$.
Every layer beyond $\ell$ consumes at least one of the $2\Delta$ excess vertices counted above.
Thus $D\le\ell+2\Delta\le b+2\Delta$.
The geometric interval sum satisfies
\[
 \sum_{0\le i\le j\le D}2^{(j-i+1)/2}
 \le \frac{\sqrt2}{(1-1/\sqrt2)^2}\,2^{D/2}
 <17\,2^{D/2}.
\]
Consequently,
\[
 c(N)<17\,2^{18\Delta+12H+18+(b+2\Delta)/2}
 \le2^{19\Delta+12H+24+b/2}.
\]
For a simple graph, no extra markings for parallel pairs are needed.
\end{proof}

\begin{theorem}[A fixed-density cycle tradeoff]
\label{thm:high-path-cycle-tradeoff}
For every $0<\varepsilon\le1/2$, there is a constant $K_\varepsilon$ such that every normalized two-pole $N$ of rank $b$ satisfies
\[
 P_{st}(N)\ge(1/2+\varepsilon)2^b
 \quad\Longrightarrow\quad
 c(N)\le K_\varepsilon 2^{b/2}.
\]
One may take $K_\varepsilon=2^{31C_\varepsilon+25}$, where $C_\varepsilon$ is the additive-stability constant.
For simple bridgeless two-poles, $2^{19C_\varepsilon+24}$ suffices.
\end{theorem}

\begin{proof}
First suppose that $N$ is bridgeless.
The additive path-stability proof gives
\[
 b-\ell\le C_\varepsilon,
 \qquad \beta(N-V(A))\le C_\varepsilon
\]
for a shortest terminal path $A$.
This application also permits a normalized multigraph.
Subdivide parallel edges to obtain a simple graph before applying the reduction.
Its suppression returns exactly $N$, and no cubic expansion is needed.
Thus the proof's residual-rank bound applies to this same $N$ and $A$.
Lemma~\ref{lem:thin-layer-cycle-count} proves the asserted estimate.

For a general normalized two-pole, every cyclic block lies on the terminal block route.
An off-route cyclic block would reduce the normalized path count to at most one half.
Every bridge therefore separates the terminals.
No bridge is incident with an original terminal, since its terminal shore would contain an off-route cycle.
Deleting all bridges leaves normalized bridgeless two-poles with distinct attachment terminals.
If two attachments coincided, the attached cyclic part would lie off the terminal route.
Every bridge shore contains a cycle, by the cubic degree count.

Write their ranks as $b_1,\ldots,b_q$.
Ranks add, terminal-path counts multiply, and cycle counts add.
Each factor's normalized path density is at least the product density, hence at least $1/2+\varepsilon$.
Apply the bridgeless bound, with constant $K_\varepsilon/2$, to every factor.
Since $b_i\ge1$ and $\sum_i b_i=b$, one has
\[
 \frac{\sum_i2^{b_i/2}}{2^{b/2}}
 \le q2^{-(q-1)/2}\le2.
\]
Summing the internal cycle counts proves the stated bound.
\end{proof}

\begin{corollary}[Simple cubic graphs with a two-edge cut]
\label{cor:simple-cubic-two-cut-error}
There is an absolute constant $K$ such that every connected bridgeless simple cubic graph $G$ of rank $r$ with a two-edge cut satisfies
\[
 c(G)\le2^{r-1}+K2^{r/2}.
\]
One may take $K=K_{1/4}$ from Theorem~\ref{thm:high-path-cycle-tradeoff}.
\end{corollary}

\begin{proof}
Deleting the cut leaves connected shores $N_1,N_2$.
The two attachment vertices in each shore are distinct.
Otherwise the shore's remaining incident edge would be a bridge of $G$.
Thus the shores are simple normalized two-poles.
Write $b_i=\beta(N_i)$, $p_i=P(N_i)$, and $c_i=c(N_i)$.
Then $b_i\ge2$, since a simple normalized two-pole cannot have rank one.
The exact decomposition is
\[
 r=b_1+b_2+1,\qquad c(G)=c_1+c_2+p_1p_2.
\]
Put $f_i=p_i/2^{b_i}$.
Call a shore high if $f_i\ge3/4$.

If both shores are high, the product term is at most $2^{r-1}$.
Theorem~\ref{thm:high-path-cycle-tradeoff} bounds the remaining sum by $K2^{r/2}$, since
$2^{b_1/2}+2^{b_2/2}\le2^{(b_1+b_2+1)/2}$.

Suppose only $N_2$ is high.
The crossing product has deficit at least $2^{b_1+b_2}/4$ from $2^{r-1}$.
Since $b_2\ge2$, this deficit is at least $2^{b_1}$.
It absorbs the trivial bound $c_1\le2^{b_1}-1$.
The remaining excess is at most $c_2$, which has the required bound.

Finally, suppose neither shore is high.
The crossing deficit is at least $7\cdot2^{b_1+b_2}/16$.
If $\max(b_1,b_2)\ge3$, then
\[
 2^{-b_1}+2^{-b_2}\le3/8<7/16.
\]
The two trivial internal cycle bounds are therefore absorbed, and $c(G)<2^{r-1}$.
The remaining case has $b_1=b_2=2$ and $r=5$.
The trivial cycle bound is already smaller than the asserted estimate.
\end{proof}

The final corollary addresses a broad decomposable subclass.
It does not prove the unrestricted half-exponent bound.
In particular, no two-edge cut is available in a cyclically three-edge-connected cubic kernel.

\begin{corollary}[All cubic two-edge cuts]
\label{cor:cubic-multigraph-two-cut-error}
There is an absolute constant $K'$ such that every connected loopless cubic multigraph with a two-edge cut satisfies
\[
 c(G)\le2^{r-1}+K'2^{r/2}.
\]
The same estimate holds for every connected cubic graph with a bridge.
\end{corollary}

\begin{proof}
First suppose that $G$ has a bridge.
Each shore contains a cycle, by the cubic degree count.
Their positive ranks $r_1,r_2$ add to $r$.
Thus
\[
 c(G)\le(2^{r_1}-1)+(2^{r_2}-1)\le2^{r-1}.
\]

Otherwise use the same two-pole decomposition as above.
Call a shore high when its normalized path density is at least $2/3$.
If both shores are high, apply Theorem~\ref{thm:high-path-cycle-tradeoff} with $\varepsilon=1/6$.
If neither is high, both ranks are at least two: a rank-one normalized two-pole is a saturated digon.
The crossing deficit is at least $(5/9)2^{b_1+b_2}$.
It absorbs both trivial internal bounds, because $2^{-b_1}+2^{-b_2}\le1/2$.

Suppose exactly one shore is high.
If its rank is at least two, its crossing deficit absorbs the low shore's internal cycles.
Indeed, $2^{b_{\rm high}}/3\ge1$.
The high shore's remaining cycles have the required half-exponent bound.

It remains to consider a high rank-one shore, which is a digon.
Let $N$ be the low shore, of rank $b$.
The exact formula gives
\[
 c(G)-2^{r-1}
 =1+c(N)+2P(N)-2^{b+1}
 \le1+c(N)-\frac23\,2^b.
\]
The asymptotic Erd\H{o}s--Simonovits bound supplies an absolute $b_0$ such that
$c(N)\le(2/3)2^b$ whenever $b\ge b_0$.
Thus the excess is at most one in this range.
For $b<b_0$, increase $K'$ to cover the trivial bound at the finitely many possible ranks.
\end{proof}

Except for the two-vertex triple-edge graph, a parallel pair in a bridgeless cubic multigraph gives a two-edge cut.
Consequently, the remaining cubic kernels for an unrestricted half-exponent bound are simple and three-edge-connected.
This reduction does not resolve that remaining class.

The exact checker reconstructs 878719 cycles on 11535 bridgeless two-poles.
Its input includes edge deletions from all 621 connected simple cubic graphs through fourteen vertices.
It also checks 94 capped-ladder and parallel-rail examples.
These checks verify the layer encoding and matching choices, rather than only the numerically loose final bound.

% Depends on the full two-pole classification and its audited local formulas.
\section{The sum bound when one terminal pair is full}
\label{sec:one-full-tripole-sum}

\begin{lemma}[One full pair]
\label{lem:one-full-tripole-sum}
Let $T$ be a connected loopless multigraph of rank $h$.
Its three distinct terminals $a,b,c$ have degree two.
Every other vertex has degree three.
If $P_{ac}(T)=2^h$, then
\[
 P_{ab}(T)+P_{bc}(T)\le2^h+6\,2^{h/2}.
\]
\end{lemma}

\begin{proof}
Suppress $b$ to an edge of a two-terminal network $N$.
Its rank remains $h$, and $P_{ac}(N)=2^h$.
The graph $N$ has no loop, since every terminal path would avoid that loop.
Lemma~\ref{lem:full-two-pole-classification} therefore applies.
Its positive-rank bridgeless pieces form a series chain.
Each piece is a capped ladder with singleton and square blocks.

First suppose $N$ consists of one such piece.
For a prefix with $p$ ladder positions, write $A_p$ for its cap count.
The empty prefix has $A_0=1$.
Adding a singleton replaces $A$ by $A+1$.
Adding a square replaces $A$ by $2A+2$.
These recurrences give
\begin{equation}\label{eq:tripole-prefix-cap}
 A_p\le3\,2^{p/2}-2\qquad(p\ge0).
\end{equation}
The singleton step uses
$3(\sqrt2-1)2^{p/2}\ge1$.
The square step preserves equality in the bound.
The same estimate holds for a suffix cap count $B_q$.

Let $p$ and $q$ count positions strictly before and after the block containing
the subdivided edge. For an interface, they count the positions on its two sides.
The subdivision formulas in Lemma~\ref{lem:one-full-three-pole-gap} give
the following table:
\[
\begin{array}{c|c|c|c}
\text{edge type}&h&P_{bc}&P_{ab}\\ \hline
\text{interface}&p+q+1&(1+A_p)2^q&(1+B_q)2^p\\
\text{singleton rung}&p+q+2&2(1+A_p)2^q&2(1+B_q)2^p\\
\text{square edge}&p+q+3&(4+3A_p)2^q&(4+3B_q)2^p
\end{array}
\]
The two parallel edges form the rank-one empty-ladder case.
Subdividing one gives $P_{ab}=P_{bc}=2$, which satisfies the conclusion.

For the first two rows, divide the path counts by $2^h$.
Their sum is bounded by $f_p+f_q$, where
\[
 f_j=\frac{1+A_j}{2^{j+1}}.
\]
Here the suffix uses its own cap count in the second summand.
Equation~\eqref{eq:tripole-prefix-cap} gives $f_j\le1/2$ for $j\ge3$.
Thus the required sum is at most $2^h$ if both $p,q\ge3$.
Otherwise reverse the piece if necessary and assume $p\le2$.
Bound $P_{bc}$ by $2^h$ using the affine path count.
In the interface row, the other term is at most
\[
 3\,2^{p+q/2}=3\,2^{(p-1)/2}2^{h/2}
 \le3\sqrt2\,2^{h/2}.
\]
In the singleton row, it is at most
\[
 6\,2^{p+q/2}=6\,2^{(p-2)/2}2^{h/2}
 \le6\,2^{h/2}.
\]

For a square edge, the normalized counts use
\[
 g_j=\frac{4+3A_j}{2^{j+3}}.
\]
Equation~\eqref{eq:tripole-prefix-cap} gives $g_j\le1/2$ for $j\ge2$.
If both $p,q\ge2$, their sum is at most one.
Otherwise assume $p\le1$ and again bound $P_{bc}$ by $2^h$.
Since $4+3B_q\le9\,2^{q/2}$, the remaining term is at most
\[
 9\,2^{p+q/2}=9\,2^{(p-3)/2}2^{h/2}
 \le\frac92\,2^{h/2}.
\]
This proves the bound for one piece.

Now suppose the subdivided edge lies inside a piece of rank $t$.
Let $L$ and $R$ be the sums of the piece ranks before and after this piece.
Write $A$ and $B$ for its paths from $b$ to its left and right endpoints.
The series decomposition gives
\[
 h=L+t+R,\qquad P_{ab}=2^L A,\qquad P_{bc}=2^R B.
\]
If $L,R\ge1$, each count is at most $2^{h-1}$.
If exactly one side has positive rank, reverse the chain and write $L=0$.
The one-piece result and $B\le2^t$ give
\[
 A+2^R B
 =(A+B)+(2^R-1)B
 \le2^{t+R}+6\,2^{t/2}
 \le2^h+6\,2^{h/2}.
\]
The case $L=R=0$ was already proved.

Finally, suppose the subdivided edge is a bridge of $N$.
Its two sides have positive ranks $L,R$ with $h=L+R$.
All terminal paths on each side are full, so
\[
 P_{ab}+P_{bc}=2^L+2^R\le2^{L+R}=2^h.
\]
This completes the proof.
\end{proof}

\begin{lemma}[One full pair leaves a small adjacent pair]
\label{lem:one-full-tripole-small-adjacent}
Under the hypotheses of Lemma~\ref{lem:one-full-tripole-sum},
\[
 \min\{P_{ab},P_{bc}\}\le3\,2^{3h/4}.
\]
\end{lemma}

\begin{proof}
Use the same full two-pole decomposition and local table.
Suppose first that the subdivided edge lies in a nonempty capped piece.
Its rank is $t=p+q+c$, where $c=1,2,3$ for the three table rows.
Equation~\eqref{eq:tripole-prefix-cap} bounds the two normalized local counts by
\[
 \frac32\,2^{-q/2},\qquad \frac32\,2^{-p/2},
\]
respectively. The square row actually has the smaller coefficient $9/8$.
Let $L,R$ be the sums of the ranks before and after this piece.
The two global normalized counts are therefore at most
\[
 \frac32\,2^{-R-q/2},\qquad \frac32\,2^{-L-p/2}.
\]
The sum of these two exponents is
\[
 L+R+\frac{p+q}{2}
 =h-\frac t2-\frac c2\ge\frac h2-\frac c2.
\]
At least one exponent is at least $h/4-c/4$.
Since $(3/2)2^{3/4}<3$, the smaller normalized count is at most
$3\,2^{-h/4}$.

If the active piece is a digon, both local counts are full.
The two global normalized counts are $2^{-R},2^{-L}$, with $L+R=h-1$.
Their minimum is at most $2^{-(h-1)/2}\le3\,2^{-h/4}$.
If the subdivided edge is a bridge, the same argument uses $L+R=h$.
Multiplication by $2^h$ proves the result.
\end{proof}

\section{Two large path counts in a three-terminal graph}
\label{sec:three-pole-high-pair-tradeoff}

A \emph{normalized three-pole} is a connected loopless multigraph with three distinct degree-two terminals.
Every other vertex has degree three.
Its rank $b$ and order $n$ satisfy $n=2b+1$.
Paths distinguish parallel edges.

\begin{theorem}[Two high pairs force a small third pair]
\label{thm:three-pole-two-high}
For every $0<\varepsilon\le1/2$, there is a constant $Q_\varepsilon$ such that every normalized three-pole $T$ satisfies
\[
 P_{ac}(T),P_{bc}(T)\ge(1/2+\varepsilon)2^b
 \quad\Longrightarrow\quad
 P_{ab}(T)\le Q_\varepsilon 2^{b/2}.
\]
One may take $Q_\varepsilon=2^{30C_\varepsilon+20}$, where $C_\varepsilon$ is the additive path-stability constant.
For a simple three-pole, $2^{18C_\varepsilon+20}$ suffices.
\end{theorem}

\begin{proof}
We first explain the use of additive stability in a normalized multigraph.
Suppress the third terminal $b$ when considering the pair $a,c$.
Its two neighbors are distinct.
Otherwise, its pendant digon would contribute rank while avoiding every $a$--$c$ path.
This would bound the normalized path count by one half.
The resulting graph is a normalized two-pole $N$ of the same rank and path count.

The additive stability proof applies directly to every shortest terminal path in $N$.
Indeed, subdividing parallel edges gives a simple graph whose normalized reduction returns $N$.
No cubic expansion is necessary.
Thus its terminal distance is at least $b-C_\varepsilon$.
The same lower bound holds in $T$.
Repeating the argument for the pair $b,c$ gives
\[
 d:=\min\{\operatorname{dist}_T(c,a),\operatorname{dist}_T(c,b)\}
 \ge b-C_\varepsilon.
\]

We will also need a shortest $c$--$a$ path $A$ in $T$ with
\[
 \beta(T-V(A))\le C_\varepsilon.
\]
Let $e$ be the suppressed edge in $N$.
If $\operatorname{dist}_T(c,a)=\operatorname{dist}_N(c,a)$, choose a shortest path in $N$ that avoids $e$.
Such a path exists, since any path using $e$ gains one edge when lifted.
Otherwise, the two distances differ by one, and every shortest path in $N$ lifts to a shortest path in $T$.
In either case, the residual-rank conclusion of additive stability applies to the chosen path in $N$.
Its lift has the same residual rank, since the extra vertex only subdivides an edge or becomes a leaf or isolated vertex.
This proves the required bound.

Assume first that $T$ is bridgeless.
It has no cut vertex, since it is subcubic.
Let $L_i$ be its distance layers from $c$.
Every layer $L_i$ with $1\le i<d$ separates $c$ from both other terminals, so $|L_i|\ge2$.
The terminal tail $\bigcup_{i\ge d}L_i$ contains both $a,b$.
Since $n=2b+1$, it follows that $d\le b$.
Put $\Delta=b-d\le C_\varepsilon$.

Mark $c$, the whole terminal tail, and every internal layer of size at least three.
The baseline counts are one vertex at $c$, two in each internal layer, and two in the terminal tail.
Their sum is $2d+1$, leaving excess $2\Delta$.
Using $w\le3(w-2)$ for $w\ge3$, the marked set has size at most $6\Delta+3$.
Every remaining layer has exactly two vertices.
It contains one vertex of $A$ and one vertex outside $A$.

In the multigraph case, mark every whole layer containing an endpoint of a parallel pair outside $A$.
There are at most $H:=\beta(T-V(A))$ such pairs, since their digons are independent.
These extra markings add at most $4H$ vertices.
A parallel pair cannot have just one endpoint on $A$, since each path vertex has at most one incidence outside $A$.
A parallel pair with both endpoints on $A$ joins consecutive path vertices, by geodesicity.

Fix a simple $a$--$b$ path $P$.
If $P$ stays in the terminal tail, its entire edge set lies in the marked set.
Otherwise, let $i<d$ be its minimum distance layer.
If that layer is not already marked, mark its two vertices.
The final marked set $B$ therefore satisfies
\[
 |B|\le6\Delta+4H+5.
\]
Fix all edges of $P$ incident with $B$, with at most $2^{3|B|}$ choices.

Every unmarked layer strictly between $i$ and $d$ carries two strands of $P$.
Indeed, both path endpoints lie beyond that layer, so every prefix cut meets $P$ evenly.
Connectedness forces both layer vertices to carry one backward and one forward path edge.
A same-layer path edge would give an odd prefix crossing.
Separate backward and forward caps would disconnect the path.

Across consecutive unmarked layers, these strands therefore form a perfect matching.
The matching is unique unless the cut is $K_{2,2}$ or contains a doubled edge of $A$.
Each exceptional cut has at most two choices after the extra markings.
Two such cuts cannot be consecutive, since their shared path vertex would have degree at least four.
Thus their combined contribution is at most $2^{(d-i)/2}$.
All remaining edges are fixed by the marked-edge data.

Summing over the minimum layer, and including paths wholly in the tail, gives
\[
 \begin{split}
 P_{ab}(T)
 &\le 2^{18\Delta+12H+15}
       \left(1+\sum_{i=0}^{d-1}2^{(d-i)/2}\right)\\
 &< 5\cdot2^{18\Delta+12H+15+d/2}
 \le 2^{30C_\varepsilon+20+b/2}.
 \end{split}
\]
For a simple graph, no parallel-edge markings are needed, so the term $12H$ is omitted.

Finally suppose that $T$ has a bridge.
Every bridge shore contains a cycle.
Indeed, a shore with $m$ vertices and $p$ original terminals has rank $(m-p+1)/2$, a positive integer.
The two high path counts force every bridge to separate $\{a,b\}$ from $c$.
No original terminal can be a bridge endpoint, since the corresponding terminal path would immediately leave a cyclic shore.

Delete one bridge, and retain the shore containing $a,b$.
It is a normalized three-pole with third terminal at the bridge attachment.
Its two normalized path densities remain at least $1/2+\varepsilon$, by factorization across the bridge.
Every $a$--$b$ path stays in this shore.
Its rank is smaller than $b$.
Induction on rank proves the same estimate, with the same constant.
\end{proof}

\begin{corollary}[A large-shore three-edge cut]
\label{cor:large-three-cut-half-exponent}
There are absolute constants $B_0,K$ such that the following holds.
Let $G$ be a connected cubic graph with a three-edge cut whose two shores are normalized three-poles of ranks at least $B_0$.
Then, writing $r=\beta(G)$,
\[
 c(G)\le2^{r-1}+K2^{r/2}.
\]
Parallel edges are allowed.
\end{corollary}

\begin{proof}
Write the shores as $T_1,T_2$, with ranks $b_1,b_2$.
Index each shore's terminal pairs by the same three pairs of cut edges.
Put
\[
 f_{ij}=\frac{P_j(T_i)}{2^{b_i}},\qquad
 f_i=(f_{i1},f_{i2},f_{i3}),\qquad c_i=c(T_i).
\]
Each cycle crosses the cut either zero or two times, so
\[
 r=b_1+b_2+2,\qquad
 c(G)=c_1+c_2+2^{b_1+b_2}\langle f_1,f_2\rangle.
\]

Call a coordinate high when it is at least $2/3$.
Put $Q=Q_{1/6}$ and choose
\[
 B_0\ge\max\{5,\lceil2\log_2(16Q)\rceil\}.
\]
The theorem shows that a shore with two high coordinates has third coordinate at most
$\delta_i:=Q2^{-b_i/2}\le1/16$.
Consequently, no shore has three high coordinates.

If both shores have two high coordinates in the same positions, then
\[
 \langle f_1,f_2\rangle\le2+\delta_1\delta_2.
\]
In every other case the dot product is at most $17/9$.
For two shores with one high coordinate each, this follows from $1+2(2/3)^2=17/9$.
If exactly one shore has two high coordinates, the bound is at most $1+2/3+(2/3)/16<17/9$.
If both have two high coordinates in different positions, the bound is at most $1+2/16<17/9$.
A shore with no high coordinate gives a smaller bound in every case.

If each shore has at least one high coordinate, suppress its unused terminal.
Theorem~\ref{thm:high-path-cycle-tradeoff} then gives
\[
 c_i\le K_{1/6}2^{b_i/2}.
\]
The aligned case has crossing excess at most $Q^2 2^{(b_1+b_2)/2}$.
The other cases have no crossing excess.
Since $b_1,b_2\ge2$, their internal cycle counts add to at most
$K_{1/6}2^{(b_1+b_2)/2}$.
This gives the claimed estimate.

If one shore has no high coordinate, the crossing deficit is at least
$2^{b_1+b_2}/9$.
The trivial internal bounds satisfy
\[
 c_1+c_2<2^{b_1}+2^{b_2}
 \le2^{b_1+b_2}/16,
\]
because $b_1,b_2\ge5$.
Thus $c(G)<2^{r-1}$ in this case.
Taking $K=Q^2+K_{1/6}$ proves the corollary.
\end{proof}

\begin{theorem}[The three-terminal path sum]
\label{thm:universal-tripole-sum}
There is an absolute constant $S$ such that every normalized three-pole $T$ of rank $b$ satisfies
\[
 P_{ab}(T)+P_{ac}(T)+P_{bc}(T)
 \le2^{b+1}+S2^{b/2}.
\]
One may take $S=\max\{6,Q_{1/32}\}$.
\end{theorem}

\begin{proof}
If one terminal pair is full, Lemma~\ref{lem:one-full-tripole-sum} proves the estimate with constant six.
Otherwise, Theorem~\ref{thm:terminal-path-saturation-gap} bounds every count by $(15/16)2^b$.
If their sum is at most $2^{b+1}$, the claim is immediate.
If their sum exceeds $2^{b+1}$, at least two counts exceed $(17/32)2^b$.
Indeed, one count at most $15/16$ and two counts at most $17/32$ have normalized sum at most two.
Apply Theorem~\ref{thm:three-pole-two-high} with $\varepsilon=1/32$.
It bounds the remaining count by $Q_{1/32}2^{b/2}$.
The other two counts are each at most $2^b$, proving the theorem.
\end{proof}

\begin{theorem}[Every cyclic cut of size at most three]
\label{thm:all-three-cut-half-exponent}
There is an absolute constant $K_3$ such that every connected loopless cubic multigraph $G$ with a cyclic edge cut of size at most three satisfies
\[
 c(G)\le2^{r-1}+K_3 2^{r/2},\qquad r=\beta(G).
\]
Here a cut is cyclic when both shores contain a cycle.
In a cubic graph, every nontrivial three-edge cut is cyclic.
\end{theorem}

\begin{proof}
A graph with a bridge has $c(G)\le2^{r-1}$.
A two-edge cut is covered by Corollary~\ref{cor:cubic-multigraph-two-cut-error}.
It remains to consider a three-edge cut in a graph with no smaller edge cut.

Each shore is connected.
Otherwise, its separate components would give two cuts of size at least three within the specified three-edge boundary.
The three attachment vertices on each shore are distinct.
If one vertex had two boundary edges, deleting that vertex from its shore would expose a two-edge cut.
Three boundary edges at one vertex would make that shore a singleton, contrary to cyclicity.
Thus the shores $T_1,T_2$ are normalized three-poles.
Their ranks $b_1,b_2$ are positive, and $r=b_1+b_2+2$.

Let $B_0$ be the constant from Corollary~\ref{cor:large-three-cut-half-exponent}.
If both ranks are at least $B_0$, that corollary applies.
Otherwise, reorder the shores so that $1\le b_1<B_0$ and $b_1\le b_2$.

Let $R_0$ be an integer such that every graph of rank at least $R_0$ has at most $(9/16)2^b$ cycles.
The established asymptotic cycle bound supplies such an integer.
Subdividing parallel edges extends that bound to loopless multigraphs.
Put
\[
 M=\max\{R_0,B_0+3\}.
\]
If $b_2<M$, then $r\le2M$.
The trivial cycle bound is absorbed by the constant $2^{M+1}$ in the asserted estimate.
Assume henceforth that $b_2\ge M$.

Suppose first that one terminal pair of $T_2$ has at least $(9/16)2^{b_2}$ paths.
Suppressing its unused terminal and applying Theorem~\ref{thm:high-path-cycle-tradeoff} gives
\[
 c(T_2)\le K_{1/16}2^{b_2/2}.
\]
Every path count in $T_1$ is at most $2^{b_1}$.
Theorem~\ref{thm:universal-tripole-sum} therefore bounds the crossing cycles by
\[
 2^{b_1}\sum_jP_j(T_2)
 \le2^{r-1}+S2^{b_1+b_2/2}.
\]
Since $b_1<B_0$, the last term is at most
$S2^{B_0/2}2^{(b_1+b_2)/2}$.
Also $c(T_1)<2^{b_1}\le2^{(b_1+b_2)/2}$.
Together, these estimates give the required bound with constant
$S2^{B_0/2}+K_{1/16}+1$.

Finally suppose that every terminal-pair count in $T_2$ is below $(9/16)2^{b_2}$.
The crossing cycles then number at most
\[
 \frac{27}{16}\,2^{b_1+b_2}.
\]
Their deficit from $2^{r-1}$ is at least
\[
 \frac5{16}2^{b_1+b_2}\ge\frac58\,2^{b_2}.
\]
The choice of $R_0$ gives $c(T_2)\le(9/16)2^{b_2}$.
The remaining deficit is at least $2^{b_2}/16$.
Since $b_2\ge B_0+3\ge b_1+4$, it absorbs $c(T_1)<2^{b_1}$.
Thus $c(G)<2^{r-1}$ in this case.

For explicit dependence, let $K_2$ be the two-edge-cut constant and let $K_{\mathrm{large}}$ be the large-shore constant.
The choice
\[
 K_3=\max\left\{
 K_2,K_{\mathrm{large}},2^{M+1},
 S2^{B_0/2}+K_{1/16}+1
 \right\}
\]
suffices.
\end{proof}

This result leaves cyclically four-edge-connected cubic kernels as the remaining unrestricted case.
It does not prove the half-exponent bound for that case.
No publication-priority claim is made here.

\section{Four-edge cuts: exact states and a forest estimate}
\label{sec:four-cut-states}

\subsection{Exact gluing formula}
Let $X,Y$ be connected four-poles with ranks $h,k$.
Their four distinct degree-two ports have matching labels $1,2,3,4$.
All other vertices have degree three.
Join equally labelled ports by four new edges to obtain $G$.
Then $r=\beta(G)=h+k+3$.

Write $p_{ij}=P_{ij}(X)$ and $q_{ij}=P_{ij}(Y)$.
Let $f_\alpha$ count the vertex-disjoint path pairs in $X$ with pairing
$\alpha\in\{12|34,13|24,14|23\}$.
Define $g_\alpha$ similarly in $Y$, and put $F=\sum_\alpha f_\alpha$,
$J=\sum_\alpha g_\alpha$.
The exact cycle decomposition is
\begin{equation}\label{eq:four-cut-exact}
 c(G)=c(X)+c(Y)+\sum_{i<j}p_{ij}q_{ij}
       +FJ-\sum_\alpha f_\alpha g_\alpha.
\end{equation}
Indeed, a crossing cycle uses either two or four cut edges.
Two cut edges give one corresponding terminal path in each shore.
Four cut edges give a path pair in each shore.
Different pairings combine into one cycle, while equal pairings give two cycles.
The half-density target for the crossing contribution is $4\,2^{h+k}$.

\subsection{At most two cycles in a graph with many terminal paths}

\begin{lemma}[Two-component cycle packings]
\label{lem:high-path-two-cycle-packings}
For each $\varepsilon>0$, there is a constant $B_\varepsilon$ with the
following property. Let $N$ be a normalized two-pole of rank $b$ with
\[
 P(N)\ge(1/2+\varepsilon)2^b.
\]
The number of nonempty unions of at most two vertex-disjoint cycles in $N$
is at most
\[
 B_\varepsilon(b+1)^4 2^{b/2}.
\]
\end{lemma}

\begin{proof}
First assume $N$ is bridgeless.
Use the shortest terminal path, distance layers, and initial marked set from
Lemma~\ref{lem:thin-layer-cycle-count}.
The additive stability theorem bounds both $\Delta$ and $H$ by
$C_\varepsilon$. The initial marked set has at most
$6\Delta+4H+2$ vertices.

Choose the minimum and maximum distance layers of each selected cycle.
There are at most $(D+1)^4$ choices for two cycles, where
$D\le b+2\Delta$ is the largest distance.
For one cycle, the corresponding bound is $(D+1)^2$.
Mark the chosen extreme layers when they were previously ordinary.
At most eight vertices are added, so the marked set $B$ satisfies
\[
 |B|\le6\Delta+4H+10.
\]
Fix all selected edges incident with $B$, at a cost of at most $2^{3|B|}$.

Every remaining ordinary layer used by one cycle lies strictly between
that cycle's extreme layers. The cycle therefore uses both vertices, with
two strands through the layer.
Two vertex-disjoint cycles cannot both use an ordinary layer.
The union thus uses each ordinary layer either through two strands or not at all.
Across consecutive occupied ordinary layers, its edges form a perfect matching.
The two-choice cuts remain nonconsecutive, as in the thin-layer lemma.
Their number is at most $(D+1)/2$.
The extreme layers, marked incidences, and matching choices determine the union.
Consequently the number of these unions is at most
\[
 \bigl((D+1)^4+(D+1)^2\bigr)
 2^{3(6\Delta+4H+10)+(D+1)/2}.
\]
Since $\Delta,H\le C_\varepsilon$, this is at most
$B'_\varepsilon(b+1)^4 2^{b/2}$ for a constant $B'_\varepsilon$.

For a general normalized two-pole, use the series decomposition from
Theorem~\ref{thm:high-path-cycle-tradeoff}.
Every positive-rank piece has normalized path density at least
$1/2+\varepsilon$.
Let their ranks be $b_1,\ldots,b_q$, with sum $b$.
Packings contained in one piece contribute at most
\[
 B'_\varepsilon(b+1)^4\sum_i2^{b_i/2}
 \le2B'_\varepsilon(b+1)^4 2^{b/2}.
\]
Packings with one cycle in each of two pieces contribute at most
\[
 K_\varepsilon^2\sum_{i<j}2^{(b_i+b_j)/2}
 \le K_\varepsilon^2b^2 2^{b/2}.
\]
Increase the constant to cover both terms.
\end{proof}

\begin{corollary}[A cross-sum forest bound]
\label{cor:four-pole-cross-sum-forest}
For each $\eta>0$, there is a constant $A_\eta$ with the following property.
Let $K$ be a connected four-pole of rank $h$ with left ports $L_1,L_2$
and right ports $R_1,R_2$.
If
\[
 \sum_{i,j=1}^2P_{L_iR_j}(K)\ge(2+\eta)2^h,
\]
then the total number of four-terminal forests in $K$ is at most
\[
 A_\eta(h+3)^4 2^{h/2}.
\]
\end{corollary}

\begin{proof}
Add a new vertex $z$ adjacent to both left ports.
Add a new vertex $w$ adjacent to both right ports.
The resulting graph $N$ is a normalized two-pole of rank $h+2$.
Its terminal-path count is the displayed cross sum, so its normalized
density is at least $1/2+\eta/4$.

Append the four new edges to any four-terminal forest of $K$.
A through pair becomes one cycle containing $z,w$.
Two same-side caps become two vertex-disjoint cycles, one through each new vertex.
This operation is injective.
Apply Lemma~\ref{lem:high-path-two-cycle-packings} and absorb the factor two
from $2^{(h+2)/2}$ into $A_\eta$.
\end{proof}

\paragraph{Status.}
The exact state identity and the complete forest estimate passed the
path-stability agent's independent audit.
These estimates supply the state-counting input for
Theorem~\ref{thm:all-four-cut-half-exponent}.
That later theorem combines them with the aggregate inequality and the
separate analysis of bounded-rank shores.

\section{Four-edge cuts with two large-rank shores}
\label{sec:four-cut-large-shores}

\begin{theorem}\label{thm:four-cut-large-shores}
There are absolute constants $H_0,K$ with the following property.
Join two connected loopless cubic four-poles of ranks $h,k\ge H_0$
by a matching of their four distinct ports.
The resulting graph has rank $r=h+k+3$ and satisfies
\[
 c(G)\le2^{r-1}+Kr^8 2^{r/2}.
\]
\end{theorem}

\begin{proof}
Use the exact gluing identity~\eqref{eq:four-cut-exact}.
For a rank-$h$ shore, normalize its six path counts and three forest counts:
\[
 a_{ij}=2^{-h}P_{ij},\qquad u_\alpha=2^{-h}f_\alpha,
 \qquad F=\sum_\alpha u_\alpha.
\]
Every $a_{ij}$ and $F$ lies in $[0,1]$.
For two normalized states $(a,u),(b,v)$, the crossing contribution divided
by $2^{h+k}$ is
\begin{equation}\label{eq:four-cut-bilinear}
 \mathcal B((a,u),(b,v))
 =\sum_{i<j}a_{ij}b_{ij}+FJ-\sum_\alpha u_\alpha v_\alpha.
\end{equation}
Define the auxiliary quantity
\[
 E(a,u)=\sum_{i<j}a_{ij}^2+F^2.
\]
Then $\mathcal B\le(E(a,u)+E(b,v))/2$.

Fix $\tau=51/100$.
Call a pair high when $a_{ij}>\tau$, and form its graph on the four ports.
Theorem~\ref{thm:three-pole-two-high}, after suppressing the fourth port,
gives an absolute constant $C$ such that
\begin{equation}\label{eq:four-cut-small-wedge}
 a_{ij},a_{jk}>\tau\quad\Longrightarrow\quad
 a_{ik}\le\rho_h:=C2^{-h/2}.
\end{equation}
If that suppression produced a loop, all paths between the other ports
would avoid its cycle and have density at most one half.
Thus this exceptional suppression cannot occur under the displayed hypothesis.

Corollary~\ref{cor:four-pole-cross-sum-forest} also supplies an absolute $C'$.
For every division of the ports into two sides, it gives
\begin{equation}\label{eq:four-cut-small-forest}
 \sum_{\mathrm{cross}}a_{ij}\ge201/100
 \quad\Longrightarrow\quad
 F\le\sigma_h:=C'(h+3)^4 2^{-h/2}.
\end{equation}
Increase $H_0$ whenever necessary below.
Both $\rho_h$ and $\sigma_h$ tend uniformly to zero for $h\ge H_0$.
In particular, the high-pair graph is triangle-free.

We classify three exceptional types of shore.
Type S has a three-edge star of high pairs, all with density one.
Type Q has a high $K_{2,2}$, all four of whose densities are one.
Type M has exactly two disjoint high pairs, both with density one.
Every other shore is called ordinary.
We claim that, for sufficiently large $H_0$,
\begin{equation}\label{eq:four-cut-energy-gap}
 E\le3.94\quad\hbox{for ordinary shores},\qquad
 E\le4.05\quad\hbox{for all shores}.
\end{equation}

Here are all cases of this finite classification.
A high $K_{2,2}$ has cross sum greater than $4\tau>2.01$.
Its two other pairs have density at most $\rho_h$, and $F\le\sigma_h$.
If one of its four high pairs is not full, the path gap bounds it by $15/16$.
Thus $E\le3+(15/16)^2+2\rho_h^2+\sigma_h^2$ in the ordinary case.
For type Q, the corresponding bound is $4+2\rho_h^2+\sigma_h^2$.

A three-edge high graph is a star or a path.
For a star, all three remaining pairs have density at most $\rho_h$.
An ordinary star gives $E\le2+(15/16)^2+3\rho_h^2+1$.
Type S gives $E\le4+3\rho_h^2$.
For a high path, use its bipartition as the division into two sides.
Its two diagonals are at most $\rho_h$, and its fourth cross pair is at most $\tau$.
If the cross sum is at least $2.01$, then
$E\le3+\tau^2+2\rho_h^2+\sigma_h^2$.
Otherwise use $a_{ij}^2\le a_{ij}$ on the cross pairs to get
$E\le2.01+2\rho_h^2+1$.

Two adjacent high pairs force their third pair below $\rho_h$.
Consequently $E\le2+3\tau^2+\rho_h^2+1$ in that case.
For two disjoint high pairs, if one is not full, then
\[
 E\le1+(15/16)^2+4\tau^2+1=3.91930625.
\]
For type M, the upper bound is $2+4\tau^2+1=4.0404$.
With at most one high pair, one has $E\le1+5\tau^2+1=3.3005$.
These cases prove~\eqref{eq:four-cut-energy-gap}.

If either shore is ordinary, the crossing contribution is at most
$3.995\,2^{h+k}$.
The trivial internal-cycle bound contributes less than $2^h+2^k$.
For large $H_0$, its normalized value $2^{-h}+2^{-k}$ is below $0.005$.
Thus such a graph has fewer than $2^{r-1}$ cycles.

We next record a restriction on type M.
Let its full matching be $\alpha=ij|kl$.
Add an edge $kl$ to the shore, obtaining a normalized two-pole $N$ of rank $h+1$.
Its terminals are $i,j$, and its counts satisfy
\[
 P_{ij}(N)=2^h\bigl(1+F-u_\alpha\bigr),\qquad c(N)\ge2^h.
\]
The first identity separates paths using the added edge from those avoiding it.
If $F-u_\alpha\ge1/100$, Theorem~\ref{thm:high-path-cycle-tradeoff}
would give $c(N)\le K_{1/200}2^{(h+1)/2}$.
This contradicts $c(N)\ge2^h$ for sufficiently large $h$.
Hence every type M shore in the present range satisfies
\begin{equation}\label{eq:four-cut-M-concentration}
 F-u_\alpha<1/100.
\end{equation}

It remains to pair exceptional types.
If two type M shores have the same full matching, their path dot product
is at most $2+4\tau^2=3.0404$.
Their off-diagonal forest term is below $2/100$ by
\eqref{eq:four-cut-M-concentration}.
Indeed, $FJ-u_\alpha v_\alpha
=F(J-v_\alpha)+v_\alpha(F-u_\alpha)<2/100$.
If their matchings differ, their path dot product is at most
$4\tau+2\tau^2=2.5602$, and their forest term is at most one.
Both cases have a fixed deficit from four.

For types M and S, the full matching and the full star share exactly one edge.
Their path dot product is at most
$1+2\tau+(1+2\tau)\rho_k$,
where $k$ is the star's rank.
The forest term is at most one, so the total stays below four by a fixed margin.
For types M and Q, the full matching has either two or zero edges in the
full $K_{2,2}$.
The path dot product is at most $2+2\tau+2\tau\rho_k$ or
$4\tau+2\rho_k$, respectively.
The forest term is at most $\sigma_k$.
These cases again have a fixed deficit.

Two type S shores with different star centers have path dot product at most
$1+2\rho_h+2\rho_k+\rho_h\rho_k$.
The forest term is at most one.
A star and a $K_{2,2}$ share exactly two full pairs.
Their path dot product is at most $2+2\rho_h+\rho_k+\rho_h\rho_k$,
with $h$ denoting the star's rank, and their forest term is at most $\sigma_k$.
Two different $K_{2,2}$ supports share exactly two edges.
Their path dot product is at most $2+2\rho_h+2\rho_k$,
and their forest term is at most $\sigma_h\sigma_k$.
All these cases have fixed deficits for large $H_0$.
The trivial internal bounds are then absorbed, as for an ordinary shore.

Only equal type S supports and equal type Q supports remain.
For equal stars, the normalized crossing contribution is at most
\[
 4+3\rho_h\rho_k.
\]
For equal $K_{2,2}$ supports, it is at most
\[
 4+2\rho_h\rho_k+\sigma_h\sigma_k.
\]
Each of these shores has a full terminal pair.
Suppressing its other two ports preserves its rank and cycle count.
The saturated two-pole cycle bound therefore gives
$c(X)\le10\,2^{h/2}$ and $c(Y)\le10\,2^{k/2}$.
The displayed crossing errors, after multiplication by $2^{h+k}$, are
at most a constant times
$(h+3)^4(k+3)^4 2^{(h+k)/2}$.
The internal counts fit the same bound.
Since $r=h+k+3$, this proves the theorem.
\end{proof}

\paragraph{Role of the rank threshold.}
The threshold $H_0$ is absolute but very large.
Theorem~\ref{thm:all-four-cut-half-exponent} completes the bounded-rank
shore cases and covers every cyclic cut of size at most four.

\section{The four-pole aggregate and every cyclic four-edge cut}
\label{sec:all-four-cuts}

For a rank-$h$ four-pole, use normalized path counts $a_{ij}$ and normalized
forest-pairing counts $u_\alpha$ as in Section~\ref{sec:four-cut-large-shores}.
Put
\[
 S=\sum_{i<j}a_{ij},\qquad F=\sum_\alpha u_\alpha.
\]
Each path count and $F$ lies in $[0,1]$.

\begin{lemma}[The added-edge constraint]\label{lem:four-pole-added-edge}
For each $0<\delta,\epsilon<1$, there is an absolute rank threshold
$H_{\delta,\epsilon}$ such that the following holds.
Assume the four-pole rank satisfies $h\ge H_{\delta,\epsilon}$.
Let $\alpha=ij|kl$ be one of the three opposite pairings.
If $a_{kl}\ge\delta$, then
\[
 a_{ij}+F-u_\alpha\le1+\epsilon.
\]
If both $a_{ij},a_{kl}\ge\delta$, then
\[
 \max(a_{ij},a_{kl})+F-u_\alpha\le1+\epsilon.
\]
\end{lemma}

\begin{proof}
Add an edge $kl$ to the four-pole.
This gives a normalized two-pole $N$ with terminals $i,j$ and rank $h+1$.
Its terminal-path and cycle counts satisfy
\[
 P_{ij}(N)=2^h(a_{ij}+F-u_\alpha),\qquad c(N)\ge a_{kl}2^h.
\]
Paths using the new edge correspond exactly to the two cross forest pairings.
Cycles using it correspond to the original $k$--$l$ paths.
If the asserted inequality failed, the normalized path density of $N$
would exceed $1/2+\epsilon/2$.
Theorem~\ref{thm:high-path-cycle-tradeoff} would imply
\[
 \delta2^h\le c(N)\le K_{\epsilon/2}2^{(h+1)/2},
\]
which is impossible for sufficiently large $h$.
Reverse the roles of the opposite pairs to obtain the second assertion.
\end{proof}

\begin{theorem}[The four-pole aggregate]\label{thm:four-pole-aggregate}
There is an absolute constant $A_4$ such that every connected loopless
cubic four-pole of rank $h$ satisfies
\[
 \sum_{i<j}P_{ij}+\sum_\alpha f_\alpha
 \le4\,2^h+A_4(h+3)^4 2^{h/2}.
\]
\end{theorem}

\begin{proof}
Suppress each unused fourth port and apply
Theorem~\ref{thm:universal-tripole-sum} to the other three.
Summing the four resulting inequalities counts each pair twice.
Consequently, for an absolute constant $C$,
\begin{equation}\label{eq:four-pole-total-paths}
 S\le4+C2^{-h/2}.
\end{equation}
If suppression creates a loop, every path between the remaining terminals
avoids its cycle and has density at most one half.
The relevant three-terminal sum bound therefore remains valid in that case.

If some cross sum is at least $2.01$, the cross-sum forest bound gives
$F\le C'(h+3)^4 2^{-h/2}$, proving the result.
Assume that every cross sum is below $2.01$.
The three cross sums add to $2S$, so $S<3.015$.
There is nothing to prove if $S+F\le4$.
Otherwise $S>3$ and $F>0.985$.

For each opposite pairing $\alpha$, let $x_\alpha$ be the sum of its two
path densities. Its complementary cross sum is $S-x_\alpha$.
Thus
\begin{equation}\label{eq:four-pole-opposite-positive}
 x_\alpha>0.99\qquad\hbox{for every }\alpha.
\end{equation}
Suppose first that no path pair is full.
The terminal-path gap bounds each density by $15/16$.
Equation~\eqref{eq:four-pole-opposite-positive} then makes every density
greater than $0.99-15/16=0.0525$.
Apply Lemma~\ref{lem:four-pole-added-edge} with $\delta=\epsilon=0.01$.
Writing $M_\alpha$ for the larger density in opposite pair $\alpha$, summation gives
\[
 \sum_\alpha M_\alpha+2F\le3.03.
\]
Since $S\le2\sum_\alpha M_\alpha$, we get $S+4F\le6.06$.
This contradicts $S>3$ and $F>0.985$.

Now suppose that $a_{12}=1$.
Lemma~\ref{lem:one-full-tripole-small-adjacent}, applied after suppressing
the unused fourth port, gives one density at most $3\,2^{-h/4}$ in each set
\[
 \{a_{13},a_{23}\},\qquad \{a_{14},a_{24}\}.
\]
For sufficiently large $h$, these two densities are at most $0.01$.
Their opposite densities exceed $0.98$ by
\eqref{eq:four-pole-opposite-positive}, and hence are full by the $15/16$ gap.
The two small densities cannot themselves be opposite, since their sum
would violate~\eqref{eq:four-pole-opposite-positive}.
Therefore the three full pairs form a star.
Theorem~\ref{thm:three-pole-two-high} bounds each of its three leaf-pair
densities by $C''2^{-h/2}$.
It follows that $S+F\le4+3C''2^{-h/2}$.

All thresholds used above are absolute.
For the finitely many smaller ranks, the trivial bound $S+F\le7$ is
absorbed by increasing $A_4$.
\end{proof}

\begin{lemma}[Low path states]\label{lem:four-pole-low-state}
For sufficiently large $h$, suppose all six normalized path densities
of a four-pole are at most $\tau=0.51$.
Then
\[
 S+F\le3.82.
\]
For any selected pairing $\gamma$, one also has
\[
 S+\frac12F+\frac12u_\gamma\le3.7.
\]
\end{lemma}

\begin{proof}
If $S\le2.75$, the first bound follows from $F\le1$.
Otherwise every path density exceeds $2.75-5\tau=0.20$.
Apply the added-edge constraint with $\delta=\epsilon=0.01$ to all three pairings.
It gives $S/2+2F\le3.03$.
Since $S\le6\tau=3.06$,
\[
 S+F\le1.515+\frac34S\le3.81<3.82.
\]

For the second bound, the case $S\le2.7$ follows from $u_\gamma\le F\le1$.
Otherwise every density exceeds $2.7-5\tau=0.15$.
Let $\alpha,\beta$ be the other pairings.
The two corresponding added-edge constraints give
\[
 M_\alpha+M_\beta+F+u_\gamma\le2.02.
\]
Consequently,
\[
 S+\frac12(F+u_\gamma)
 \le S-\frac12(M_\alpha+M_\beta)+1.01.
\]
The two coordinates in each of $\alpha,\beta$ contribute at most
$3\tau/2$ after this subtraction.
The coordinates in $\gamma$ contribute at most $2\tau$.
The last display is therefore at most $5\tau+1.01=3.56<3.7$.
\end{proof}

\begin{theorem}[Every cyclic cut of size at most four]
\label{thm:all-four-cut-half-exponent}
There is an absolute constant $K_4$ such that every connected loopless
cubic multigraph with a cyclic edge cut of size at most four satisfies
\[
 c(G)\le2^{r-1}+K_4r^8 2^{r/2},\qquad r=\beta(G).
\]
A cut is cyclic when both shores contain a cycle.
\end{theorem}

\begin{proof}
Cyclic cuts of size at most three are covered by
Theorem~\ref{thm:all-three-cut-half-exponent}.
Assume none exists, and consider a cyclic cut of size four.
The graph has no bridge or two-edge cut, since those are cyclic in a cubic graph.
Thus each component after deleting the four cut edges has boundary at least three.
There are exactly two components: three would require total boundary at least nine,
whereas the four cut edges supply only eight.
Write these connected cyclic shores as $X,Y$.

Their attachment vertices are distinct.
Indeed, a vertex with two cut edges has one internal edge and is a leaf in its shore.
Move this vertex to the other shore.
The original shore remains connected and cyclic, while the other remains connected
and gains a cycle. The new boundary has size three, a contradiction.
A vertex with three cut edges would make its connected shore a singleton,
contrary to cyclicity.
Thus $X,Y$ are normalized four-poles.
Their positive ranks $h,k$ satisfy $r=h+k+3$.

If $h,k\ge H_0$, apply Theorem~\ref{thm:four-cut-large-shores}.
Otherwise reorder so that $1\le h<H_0$ and $h\le k$.
If $k$ is bounded by any fixed threshold needed below, all ranks are bounded.
Increase $K_4$ to cover those cases by the trivial cycle bound.
Assume henceforth that $k$ exceeds all these absolute thresholds.

Suppose some terminal pair of $Y$ has at least $0.51\,2^k$ paths.
Suppress the other two ports and apply
Theorem~\ref{thm:high-path-cycle-tradeoff} to obtain
$c(Y)\le K_{1/100}2^{k/2}$.
Any loop created by suppression would force path density at most one half,
so it cannot occur in this case.
Every path count of $X$ is at most $2^h$.
Every forest coefficient $\sum_\alpha f_\alpha-f_\beta$ is also at most $2^h$.
The exact gluing formula and the aggregate theorem therefore give
\[
 c(G)\le c(X)+c(Y)+2^h
 \left(4\,2^k+A_4(k+3)^4 2^{k/2}\right).
\]
The first crossing term is $2^{r-1}$.
Since $h<H_0$, all remaining terms are
$O(r^4 2^{r/2})$, with an absolute implied constant.

It remains to suppose that every path density of $Y$ is at most $0.51$.
The established asymptotic cycle bound gives
\[
 c(Y)\le\frac9{16}2^k
\]
for sufficiently large $k$.
This also applies to multigraphs after uniform subdivision.
If $h\ge2$, the first bound in Lemma~\ref{lem:four-pole-low-state}
limits crossing cycles to $3.82\,2^{h+k}$.
Their deficit from $2^{r-1}$ is at least
\[
 0.18\,2^{h+k}\ge0.72\,2^k.
\]
It absorbs $c(Y)\le(9/16)2^k$ and $c(X)<2^h$ for large $k$,
because $h<H_0$ is bounded.
Thus $c(G)<2^{r-1}$ in this case.

Finally suppose $h=1$.
The degree sum gives four vertices in $X$, all of degree two.
Since $X$ is connected and loopless, it is a four-cycle.
Its six path counts are all two.
Its normalized forest vector has entries $1/2,1/2,0$.
Let $\gamma$ be the pairing with zero entry.
The exact normalized crossing contribution is
\[
 S+\frac12F+\frac12u_\gamma,
\]
where the state belongs to $Y$.
Lemma~\ref{lem:four-pole-low-state} bounds this expression by $3.7$.
The crossing deficit is therefore at least $0.6\,2^k$.
This absorbs $c(Y)\le(9/16)2^k$ and $c(X)=1$ for sufficiently large $k$.
The theorem follows after increasing $K_4$ over the bounded cases.
\end{proof}

\paragraph{Scope.}
Cyclicity is essential to the cut hypothesis.
An acyclic two-vertex shore has a four-edge boundary in every cubic graph.
The theorem does not prove an unrestricted half-exponent bound.
After the earlier reductions, the remaining cubic kernels are cyclically
five-edge-connected.

% Independent mathematical audits passed: root, path_stability, and random_hosts.
\section{A linear factor for two-cycle packings}
\label{sec:linear-two-cycle-packings}

This refinement improves the fourth-degree polynomial in Lemma~\ref{lem:high-path-two-cycle-packings} to a linear factor.
It uses the same marked layers and matching choices.
The improvement comes from summing the possible cycle intervals with their geometric weights.

\begin{lemma}[Two intervals with bounded overlap]
\label{lem:bounded-overlap-interval-sum}
Fix a nonnegative integer $B$.
There is a constant $c_B$ such that, for every integer $D\ge0$,
\[
 \sum_{\substack{I,J\text{ nonempty integer intervals in }[0,D]\\
                  |I\cap J|\le B}}
 2^{(|I|+|J|)/2}
 \le c_B(D+1)2^{D/2}.
\]
The intervals in this sum are ordered.
\end{lemma}

\begin{proof}
Put $u=\sqrt2$.
The one-interval sum satisfies
\[
 \sum_{0\le i\le j\le D}u^{j-i+1}
 \le\frac{u}{(1-u^{-1})^2}u^D<17u^D.
\]
Order the left endpoints, at a cost of at most two.
There are three cases.

First suppose that the intervals are disjoint, with the first before the second.
Let their outer endpoints be $i,l$, and let $g$ count the intervening unused positions.
Their total length is $l-i+1-g$.
For fixed $i,l,g$, there are at most $D+1$ choices of the split position.
Summing the geometric factor $u^{-g}$ gives a bound
\[
 \frac{17}{1-u^{-1}}(D+1)u^D.
\]

Next suppose that the intervals overlap without one containing the other.
Let $q\le B$ be their overlap length.
Their total length is the outer span plus $q$.
For fixed outer endpoints and $q$, there are at most $D+1$ meeting positions.
The resulting bound is
\[
 17\left(\sum_{q=1}^{B}u^q\right)(D+1)u^D.
\]
Finally, if one interval contains the other, the inner interval has length at most $B$.
For each outer interval and inner length, there are at most $D+1$ placements.
The same bound applies.
Adding the three cases proves the lemma.
\end{proof}

\begin{theorem}[Linear two-cycle packing bound]
\label{thm:linear-high-path-two-cycle-packings}
For every fixed $0<\varepsilon\le1/2$, there is a constant $B_\varepsilon$ such that the following holds.
Let $N$ be a normalized two-pole of rank $b$ with
\[
 P(N)\ge(1/2+\varepsilon)2^b.
\]
The number of nonempty unions of at most two vertex-disjoint cycles in $N$ is at most
\[
 B_\varepsilon(b+1)2^{b/2}.
\]
\end{theorem}

\begin{proof}
First suppose that $N$ is bridgeless.
Use the shortest terminal path and initial whole-layer marking from Lemma~\ref{lem:thin-layer-cycle-count}.
Additive stability bounds both $\Delta$ and $H$ by $C_\varepsilon$.
Thus the initial marked set has at most $10C_\varepsilon+2$ vertices and meets at most that many layers.
Write $D$ for the maximum distance layer, so $D\le b+2C_\varepsilon$.

Each cycle has a nonempty interval of distance layers.
Consider two vertex-disjoint cycles with intervals $I,J$.
Any layer strictly inside both intervals must contain at least two vertices from each cycle.
An initially ordinary layer has only two vertices, so it cannot have this property.
At most four layers are interval endpoints.
Consequently,
\[
 |I\cap J|\le \lceil10C_\varepsilon+2\rceil+4=:B.
\]

Fix the two intervals, and mark all their endpoint layers.
This adds at most eight vertices to the initial marked set.
Fix all selected edges incident with marked vertices, at a cost bounded by a constant depending only on $\varepsilon$.
Every remaining ordinary layer is unused or carries the two through strands of exactly one cycle.
The intervals determine which possibility applies.
Across consecutive occupied ordinary layers, the selected edges form a perfect matching.
Two-choice cuts cannot be consecutive, by the same cubic degree argument as before.
Their number is at most $(|I|+|J|)/2$.
The interval data, marked incidences, and matching choices determine the union of cycles.

Lemma~\ref{lem:bounded-overlap-interval-sum} now bounds two-cycle unions by
\[
 O_\varepsilon((D+1)2^{D/2})
 =O_\varepsilon((b+1)2^{b/2}).
\]
The one-cycle contribution is $O_\varepsilon(2^{b/2})$ by Theorem~\ref{thm:high-path-cycle-tradeoff}.
This proves the bridgeless case.

For a general normalized two-pole, use its series decomposition into bridgeless pieces.
Let their positive ranks be $b_1,\ldots,b_q$, with sum $b$.
Every piece retains the required normalized path density.
Packings contained in one piece contribute at most
\[
 B'_\varepsilon\sum_i(b_i+1)2^{b_i/2}
 \le B'_\varepsilon(b+2)2^{b/2}.
\]
Here we used $\sum_i2^{b_i/2}\le2\,2^{b/2}$.

Packings with one cycle in each of two pieces contribute at most
\[
 K_\varepsilon^2\sum_{i<j}2^{(b_i+b_j)/2}.
\]
For $q\ge2$,
\[
 \frac{\sum_{i<j}2^{(b_i+b_j)/2}}{2^{b/2}}
 \le\binom q2\,2^{-(q-2)/2}\le4.
\]
The last bound follows by checking $2\le q\le6$ and then taking successive ratios.
Increasing $B_\varepsilon$ absorbs both contributions.
\end{proof}

\begin{corollary}[Linear cross-sum forest bound]
\label{cor:linear-four-pole-cross-sum-forest}
For each fixed $\eta>0$, there is a constant $A_\eta$ such that every connected four-pole $K$ of rank $h$ satisfies
\[
 \sum_{i,j=1}^2P_{L_iR_j}(K)\ge(2+\eta)2^h
 \quad\Longrightarrow\quad
 F(K)\le A_\eta(h+3)2^{h/2}.
\]
Here $F(K)$ counts all vertex-disjoint path pairs connecting its four ports.
\end{corollary}

\begin{proof}
Cap the two sides by new terminals, as in Corollary~\ref{cor:four-pole-cross-sum-forest}.
The resulting two-pole has rank $h+2$ and normalized path density at least $1/2+\eta/4$.
Each four-terminal forest maps injectively to one cycle or two vertex-disjoint cycles in this two-pole.
Apply Theorem~\ref{thm:linear-high-path-two-cycle-packings}.
\end{proof}

The linear factor in the packing theorem is necessary.
Take a capped ladder consisting of $m$ square blocks.
Its terminal paths are saturated and its rank is $b=2m+1$.
Count two disjoint cycles, one through each cap terminal.
If the left cycle turns in square $i$ and the right cycle turns in square $j>i$, their number is $2^{m+i-j+1}$.
Summing these choices gives
\[
 \sum_{1\le i<j\le m}2^{m+i-j+1}
 =(m-2)2^{m+1}+4
 =\Theta(b2^{b/2}).
\]
The turns must occur in different squares because two same-square caps cannot be vertex-disjoint.
Thus a bound $O_\varepsilon(2^{b/2})$ without a rank factor would fail even at full terminal-path density.

\section{A path obstruction from the Horak--Kim construction}
\label{sec:second-order-constructions}

The proposed additive path bound needs a polynomial factor.
The following family gives a linear obstruction at terminal distance three.
It uses the graphs $MC_{2q}$ of Horak and Kim~\cite{HK}.
Their total cycle count already gives the known lower bound
$\psi(r)-2^{r-1}=\Omega(r2^{r/2})$.
The marked-edge count below gives a separate consequence for terminal paths.

\begin{proposition}\label{prop:hk-short-path-obstruction}
For every integer $q\ge2$, there is a simple connected graph $N_q$
with terminals $u,v$, rank $b=2q$, and $\operatorname{dist}_{N_q}(u,v)=3$,
such that
\[
 P_{uv}(N_q)=2^{2q-1}+\left(q-\frac34\right)2^q+1.
\]
Consequently, no absolute constant $C$ satisfies
\[
 P_{uv}(N)\le 2^{b-1}+C2^{(b+\ell)/2}
\]
for every connected graph of rank $b$ with terminal distance $\ell$.
\end{proposition}

\begin{proof}
Use indices modulo $q$.
Block $i$ has vertices $a_i,b_i,c_i,d_i$ and all four edges between
$\{a_i,b_i\}$ and $\{c_i,d_i\}$.
Add the interface edges $c_i a_{i+1}$ and $d_i b_{i+1}$.
The resulting graph $G_q$ is simple and cubic.
It has $4q$ vertices, $6q$ edges, and rank $2q+1$.
Set $e=a_0c_0$ and $N_q=G_q-e$.

Cut parity gives two cases for every cycle of $G_q$.
It uses one edge at every interface, or it uses zero or two at every interface.
In the first case, choose one edge independently at every interface.
Each square then has two paths between the selected endpoints.
There are therefore $2^q2^q=4^q$ cycles in this case.
Exactly half contain $e$.
Indeed, four of the eight local endpoint-and-path choices contain any fixed square edge.
The contribution is $2^{2q-1}$.

Now consider the second case.
If no interface is used, only the square in block zero contains $e$.
This contributes one cycle.
If some interfaces are used, but not all, they form a cyclic interval.
An interval of $d$ interfaces occupies $d+1$ blocks, where $1\le d\le q-1$.
Each end block has two possible caps.
Each inner block has two possible matchings.
There are $2^{d+1}$ cycles for each such interval.
Exactly $d+1$ intervals contain block zero.
For each of these intervals, half the cycles contain $e$.
Their total contribution is
\[
 \sum_{d=1}^{q-1}(d+1)2^d=(q-1)2^q.
\]

Finally, suppose that every interface is used.
Each square supplies one of its two matchings.
The result is one cycle precisely when the matching product is the transposition.
There are $2^{q-1}$ such choices.
Exactly half contain $e$, so they contribute $2^{q-2}$.
Adding the contributions gives
\[
 c_e(G_q)=2^{2q-1}+1+(q-1)2^q+2^{q-2}.
\]

Deleting $e$ gives a bijection from these cycles to simple terminal paths.
The remaining three edges of its square give a path of length three.
The graph is bipartite, so the terminal distance is three.
Its rank is $2q$.
The normalized excess equals
\[
 \frac{P_{uv}(N_q)-2^{b-1}}{2^{(b+3)/2}}
 =\frac{q-3/4+2^{-q}}{2^{3/2}},
\]
which tends to infinity.
\end{proof}

This obstruction does not refute an error bound $O(r2^{r/2})$ for total cycles.
It shows that a direct path bound at the half exponent needs additional rank dependence.
The independent script \texttt{check\_second\_order\_constructions.py}
counts terminal paths for $2\le q\le9$ and verifies the formula.
These finite checks support the symbolic proof; they do not establish optimality.

% Independent mathematical audit passed: random_hosts. Dependencies separately audited.
\section{The sharp order for cyclic cuts of size at most four}
\label{sec:four-cut-sharp-order}

The linear two-cycle packing bound improves both four-cut estimates.
For aligned saturated shores, the largest forest terms cancel in the gluing identity.
These two improvements give the same order as the cyclic-square lower bound.

\begin{corollary}[A linear four-pole aggregate error]
\label{cor:linear-four-pole-aggregate}
There is an absolute constant $A'_4$ such that every connected loopless cubic four-pole of rank $h$ satisfies
\[
 \sum_{i<j}P_{ij}+\sum_\alpha f_\alpha
 \le4\,2^h+A'_4(h+3)2^{h/2}.
\]
\end{corollary}

\begin{proof}
Repeat the proof of Theorem~\ref{thm:four-pole-aggregate}.
Its only fourth-degree factor comes from the cross-sum forest estimate.
Replace that estimate by Corollary~\ref{cor:linear-four-pole-cross-sum-forest}.
The low-cross-sum full-star case still has error $O(2^{h/2})$.
All fixed rank thresholds and the finitely many smaller ranks are absorbed into $A'_4$.
\end{proof}

\begin{theorem}[Sharp-order four-cut error]
\label{thm:sharp-order-four-cut-error}
There is an absolute constant $K$ such that every connected loopless cubic multigraph of rank $r$ with a cyclic edge cut of size at most four satisfies
\[
 c(G)\le2^{r-1}+Kr2^{r/2}.
\]
The factor $r$ cannot be replaced by a function that is $o(r)$ throughout this class.
\end{theorem}

\begin{proof}
Cuts of size at most three already have error $O(2^{r/2})$ by Theorem~\ref{thm:all-three-cut-half-exponent}.
For a cyclic four-edge cut, use the connected four-pole shores from Theorem~\ref{thm:all-four-cut-half-exponent}.
Write their ranks as $h,k$, so $r=h+k+3$.

First suppose both ranks exceed the fixed threshold in Theorem~\ref{thm:four-cut-large-shores}.
Increase that threshold if needed for the constant in the linear forest estimate.
Its ordinary and nonaligned exceptional cases have fixed crossing deficits.
Those deficits still absorb all internal cycles.
The same profile classification applies when the forest estimate is replaced by its linear refinement.
Only equal full-star supports and equal full-$K_{2,2}$ supports need a smaller error bound.

For equal stars, the earlier bound is
\[
 \mathcal B\le4+3\rho_h\rho_k,
 \qquad \rho_j=C2^{-j/2}.
\]
Thus their crossing error is $O(2^{(h+k)/2})$.
Both shores have a full terminal pair, so their internal cycle counts have the same half-exponent bound.

Now suppose the two full-$K_{2,2}$ supports agree.
They have the same distinguished same-side cap pairing $\alpha$.
For the rank-$h$ shore, put
\[
 C_h=f_\alpha,\qquad D_h=\sum_{\beta\ne\alpha}f_\beta.
\]
Thus $D_h$ counts through pairs and $C_h$ counts pairs of disjoint same-side caps.
Define $C_k,D_k$ in the other shore.
The audited saturated-block bounds give
\[
 D_h\le2^{(h+1)/2},\qquad
 C_h\le3(h+1)2^{(h+1)/2},
\]
and the corresponding inequalities at rank $k$.

The exact four-cut gluing identity subtracts products with equal pairings.
In particular, its $C_hC_k$ term cancels.
The remaining four-edge contribution is at most
\[
 C_hD_k+D_hC_k+D_hD_k
 \le[3(h+k+2)+1]2^{(h+k+2)/2}.
\]
The two nonfull path-coordinate products contribute only $O(2^{(h+k)/2})$.
Internal cycles again contribute $O(2^{h/2}+2^{k/2})$.
This proves the required $O(r2^{r/2})$ bound when both shores have large rank.

Suppose instead that one shore has bounded rank.
The proof of Theorem~\ref{thm:all-four-cut-half-exponent} already places every low-path-density large shore below the one-half bound.
Its only remaining case uses the universal four-pole aggregate and a high path count in the large shore.
Replace the aggregate there by Corollary~\ref{cor:linear-four-pole-aggregate}.
The resulting crossing error is $O((k+3)2^{k/2})$, since the other rank is bounded.
The large shore's internal cycles are $O(2^{k/2})$ by the fixed-density tradeoff.
All bounded ranks are absorbed into $K$.
This proves the upper bound.

For the lower bound, take the cyclic-square graph $G_q$ from Proposition~\ref{prop:hk-short-path-obstruction}.
It has rank $r=2q+1$ and a cyclic four-edge cut around each square.
There are $4^q$ cycles using one edge at every interface.
There are $q$ cycles using no interface.
A proper interval of $d$ interfaces gives $q2^{d+1}$ cycles, for $1\le d\le q-1$.
Using both edges at every interface gives $2^{q-1}$ cycles.
Therefore
\[
 \begin{split}
 c(G_q)
 &=4^q+q+q\sum_{d=1}^{q-1}2^{d+1}+2^{q-1}\\
 &=2^{2q}+\left(2q+\frac12\right)2^q-3q.
 \end{split}
\]
Its excess is $\Theta(q2^q)=\Theta(r2^{r/2})$.
\end{proof}

This theorem addresses a structural subclass of cubic kernels.
It does not prove the unrestricted sharp-order error.
After the established reductions, that remaining question concerns cyclically five-edge-connected cubic kernels.

The independent script \texttt{check\_four\_cut\_sharp\_order.py} verifies the total cycle formula for $2\le q\le7$ by binary cycle-space enumeration.
It also verifies the disjoint-cap packing formula for one through six square blocks.
Its results are saved in \texttt{four-cut-sharp-order-checks.json}.

\section{Five-edge cuts: exact identities and partial restrictions}
\label{sec:five-cut-partial}

This section records exact identities and initial restrictions for the five-edge-cut analysis.
The later universal functional completes the sharp-order cycle bound.

Let $X,Y$ be connected normalized five-poles of ranks $h,k$.
Join corresponding ports by five edges, to obtain $G$ of rank $r=h+k+4$.
For each pair $e=ij$, write $P_e,Q_e$ for the terminal-path counts.
For each four-element port set $S$, and each of its three pairings $\alpha$, write
$f_{S,\alpha},g_{S,\alpha}$ for the counts of two vertex-disjoint paths with that pairing.
Put $F_S=\sum_\alpha f_{S,\alpha}$ and $J_S=\sum_\alpha g_{S,\alpha}$.
The fifth port has even degree in each counted forest.

\begin{proposition}[Five-cut gluing]
\label{prop:five-cut-gluing}
One has the exact identity
\[
 c(G)=c(X)+c(Y)+\sum_{|e|=2}P_eQ_e
 +\sum_{|S|=4}\left(F_SJ_S-\sum_\alpha f_{S,\alpha}g_{S,\alpha}\right).
\]
\end{proposition}
\begin{proof}
A cycle uses zero, two, or four edges of the cut.
The zero-edge and two-edge cases give the first three terms.
For a fixed four-edge set, each shore supplies a pairing forest.
Equal pairings give two cycles, while distinct pairings give one cycle.
Thus the four-edge contribution has the displayed form.
\end{proof}

Normalize all shore counts by $2^h$ or $2^k$ as appropriate.
Write $a_e=2^{-h}P_e$, $u_{S,\alpha}=2^{-h}f_{S,\alpha}$, and $U_S=\sum_\alpha u_{S,\alpha}$.
Every coordinate $a_e$ and every sum $U_S$ lies in $[0,1]$.
Define the shore energy
\[
 E(X)=\sum_e a_e^2+\sum_S U_S^2.
\]
The normalized crossing contribution is at most $(E(X)+E(Y))/2$.
The leading one-half threshold for that contribution is $8$.
This energy estimate drops the nonnegative equal-pairing subtraction, so it can lose important information.

\begin{proposition}[Capping three ports]
\label{prop:five-pole-three-cap}
Choose three ports $A$ and let $d,e$ be the other two.
Add a new vertex adjacent to the three ports of $A$.
The resulting two-pole $N$ has rank $h+2$, and
\[
 P_{de}(N)=P_{de}(X)+\sum_{\{a,b\}\subset A}
 \left(F_{\{d,e,a,b\}}-f_{\{d,e,a,b\},de|ab}\right),
\]
\[
 c(N)=c(X)+\sum_{\{a,b\}\subset A}P_{ab}(X).
\]
For any fixed $\delta,\eta>0$, if some $a_{ab}\ge\delta$ with $a,b\in A$, then, for all sufficiently large $h$,
\[
 a_{de}+\sum_{\{a,b\}\subset A}
 \left(U_{\{d,e,a,b\}}-u_{\{d,e,a,b\},de|ab}\right)\le2+\eta.
\]
\end{proposition}
\begin{proof}
A terminal path avoids the new vertex or traverses two of its incident edges.
In the latter case, deleting the new vertex gives either of the two pairings other than $de|ab$.
A cycle through the new vertex gives an $a$--$b$ path in $X$.
This proves the identities.
If the last inequality failed, $N$ would have normalized terminal-path density greater than $1/2+\eta/4$.
The fixed-density path--cycle tradeoff would give $c(N)=O_\eta(2^{h/2})$.
This contradicts $c(N)\ge\delta2^h$ for large $h$.
\end{proof}

\begin{proposition}[Adding one edge]
\label{prop:five-pole-added-edge}
Choose two ports $i,j$ and let $R$ be the other three.
Adding edge $ij$ gives a tripole $T$ of rank $h+1$.
For each pair $e\subset R$,
\[
 2^{-h}P_e(T)=a_e+U_{\{i,j\}\cup e}-u_{\{i,j\}\cup e,ij|e}.
\]
Consequently,
\[
 \sum_{e\subset R}\left(a_e+U_{\{i,j\}\cup e}-u_{\{i,j\}\cup e,ij|e}\right)
 \le4+O(2^{-h/2}).
\]
If two of the three summands exceed $1+\eta$ for fixed $\eta>0$, the third is $O_\eta(2^{-h/2})$.
\end{proposition}
\begin{proof}
Separate paths that use the added edge from those that avoid it.
Apply the universal tripole path-sum estimate at rank $h+1$.
For the last assertion, divide all three counts by $2^{h+1}$ and apply the two-high-pairs tripole theorem.
\end{proof}

\begin{proposition}[A gap when all path coordinates are positive]
\label{prop:five-pole-positive-energy-gap}
For all sufficiently large $h$, a normalized five-pole with $a_e\ge1/100$ for all ten pairs satisfies
\[
 E(X)\le7.9483<8.
\]
\end{proposition}
\begin{proof}
Call a pair high if $a_e>51/100$.
The two-high-pairs tripole theorem, after suppressing the unused ports, implies that a wedge of high pairs has its third pair $O(2^{-h/2})$.
Thus the high-pair graph is a matching for large $h$, and it has at most two edges.

Apply the four-pole added-edge constraint to each four-port set, with $\delta=1/100$ and $\eta=1/1000$.
Unused degree-two ports may be suppressed.
If suppression creates a loop, that loop is at a nonterminal vertex.
Every path in the added-edge two-pole avoids it.
Its normalized path count is then at most one half, which gives the same constraint directly.
Summing its six oriented constraints gives
\[
 \sum_{e\subset S}a_e+4U_S\le6(1+\eta).
\]
Summing over the five four-port sets counts each pair three times, so
\[
 3\sum_ea_e+4\sum_SU_S\le30(1+\eta).
\]
Since $U_S^2\le U_S$,
\[
 E(X)\le\frac{15}{2}(1+\eta)+\sum_e(a_e^2-3a_e/4).
\]
The summand is at most $1/4$ on a high pair.
For a nonhigh pair it is at most $-0.0074$, by convexity on $[0.01,0.51]$.
There are at most two high pairs.
Therefore
\[
 E(X)\le7.5(1.001)+2(0.25)-8(0.0074)=7.9483.
\]
\end{proof}

The independent script \texttt{check\_five\_cut\_states.py} verifies 1,920 complete five-cut gluings and 40 three-port cappings on four small shores.
The checks include a shore with a bridge.
The files \texttt{explore\_five\_cut\_profiles.py} and \texttt{five-cut-profile-lp-exploration.json} contain numerical relaxations only; they are not proof certificates.
The universal functional proved later resolves both sparse profiles and bounded shores for the cycle bound.

\section{A high path density forces a small cyclic cut}
\label{sec:high-density-small-cut}

The saturated-pair conclusion extends to every fixed path density above one half.
The proof uses additive path stability and a local property of distance layers.
No classification of saturated two-poles is needed.

\begin{lemma}[Four consecutive layers]
\label{lem:four-thin-layers-short-cycle}
Let $N$ be a bridgeless normalized two-pole, with distance layers from one terminal.
Suppose four consecutive internal layers each have exactly two vertices.
Then their union contains a cycle of length at most four.
\end{lemma}

\begin{proof}
Every vertex in these layers has degree three and a neighbor in the preceding layer.
Edges join vertices in the same layer or consecutive layers.
Each of the three cuts between these layers has size at least two, since $N$ is bridgeless.
Each cut has size at most four, because its earlier layer has two predecessor incidences.

A parallel pair inside these four layers already gives a digon.
Assume no such pair occurs.
A cut of size four then induces $K_{2,2}$ and gives a four-cycle.
Assume instead that all three cut sizes belong to $\{2,3\}$.
Write them as $e_1,e_2,e_3$.
For either middle layer, let $r_i\in\{0,1\}$ indicate its internal edge.
Its degree sum gives
\[
 e_1+e_2+2r_2=6,\qquad e_2+e_3+2r_3=6.
\]
These identities force $e_1=e_2=e_3$.

If all sizes are two, both middle layers have an internal edge.
Each vertex in either middle layer has one preceding and one following edge.
The cut between them is therefore a perfect matching.
Together with their two internal edges, it forms a four-cycle.

If all sizes are three, both middle layers have no internal edge.
Use the first three layers.
In the first cut, one middle-layer vertex $u$ has two predecessors and the other vertex $v$ has one.
Call the predecessor of $v$ in the first layer $x$.
Since $u$ has both possible predecessors, $xu$ is also an edge.
The cubic degree condition makes $u$ have one successor and $v$ have two successors.
Call the successor of $u$ in the third layer $y$.
Then $vy$ is also an edge, so $xuyvx$ is a four-cycle.
\end{proof}

\begin{theorem}[A fixed-density pair inside a fixed-order pole]
\label{thm:high-density-pair-small-cyclic-cut}
Fix $k\ge2$ and $0<\varepsilon\le1/2$.
Let $C_\varepsilon$ be the additive path-stability constant, and put
\[
 B_{k,\varepsilon}
 =(k+1)\bigl(9\lceil C_\varepsilon\rceil+8k+9\bigr).
\]
Let $X$ be a connected normalized $k$-pole of rank $h$.
If one terminal pair satisfies
\[
 P_{st}(X)\ge(1/2+\varepsilon)2^h,
 \qquad h\ge B_{k,\varepsilon},
\]
then every cubic completion of $X$ has a cyclic edge cut of size at most four.
\end{theorem}

\begin{proof}
Suppress every port except $s,t$, obtaining a normalized two-pole $N$ of the same rank and terminal-path count.
It has no loop, since an independent loop coordinate would force path density at most one half.
Every shore of a bridge contains a cycle, by the degree count for a normalized two-pole.
Therefore every bridge separates $s$ from $t$.
Otherwise, all terminal paths would avoid a cyclic shore and have density at most one half.
The same argument excludes original terminals at bridge endpoints and cyclic parts outside the terminal route.

Thus the positive-rank bridgeless pieces form a series chain, with ranks $b_1,\ldots,b_q$ satisfying $\sum b_i=h$.
Their attachment vertices are distinct, since coincident attachments would leave a cyclic part outside every terminal path.
Each piece is a normalized two-pole.
Path counts multiply across the chain.
Each piece consequently has normalized terminal-path density at least $1/2+\varepsilon$.

Suppose first that $q\ge k+2$.
Exclude the first and last pieces and every piece containing an internal subdivision by a restored port.
At most $k$ pieces are excluded, so at least two remain.
Each remaining whole piece contains a cycle and has boundary at most two in every completion.
Subdivisions of its boundary edges do not enlarge that boundary.
One piece supplies a cyclic cut, and the other supplies a cycle in its complement.

Assume henceforth that $q\le k+1$.
Some bridgeless piece $M$ has rank
\[
 b\ge\frac h{k+1}\ge9\lceil C_\varepsilon\rceil+8k+9.
\]
Write $C=\lceil C_\varepsilon\rceil$.
Additive path stability gives terminal distance $\ell\ge b-C$ in $M$.
The graph $M$ has $2b$ vertices.
It is two-connected, because a bridgeless subcubic graph cannot have a cut vertex.
Hence each internal distance layer $L_i$, $1\le i<\ell$, has at least two vertices.
There are at most $2(b-\ell)\le2C$ such layers with more than two vertices.

Call those larger layers bad.
Also call a layer bad if it contains an endpoint of an edge subdivided by a restored port.
There are at most $2(k-2)$ additional bad layers.
Partition $L_1,\ldots,L_{\ell-1}$ into consecutive groups of four, discarding the final incomplete group.
The number of groups without a bad layer is at least
\[
 \left\lfloor\frac{\ell-1}{4}\right\rfloor-2C-2(k-2)
 \ge\left\lfloor\frac{b-C-1}{4}\right\rfloor-2C-2k+4
 \ge2.
\]
The displayed rank threshold has ample slack for the last inequality.

Choose two such groups.
Lemma~\ref{lem:four-thin-layers-short-cycle} supplies one cycle of length at most four in each group.
The cycles are vertex-disjoint.
Neither cycle uses an edge subdivided by a restored port.
Both therefore survive unchanged in $X$ and in every completion.

In a cubic graph, the vertex set of a cycle of length at most four has boundary at most four.
Choose one of these two cycle vertex sets.
The other cycle lies in its complement, so its boundary is a cyclic edge cut.
This proves the theorem.
\end{proof}

\begin{corollary}[Path density inside a cyclically five-connected completion]
\label{cor:cyclic-five-path-density}
Fix $k\ge2$ and $\varepsilon>0$.
For all sufficiently large ranks, every normalized $k$-pole inside a cubic graph without a cyclic cut of size at most four satisfies
\[
 P_{st}<(1/2+\varepsilon)2^h
\]
for every pair of its ports.
The rank threshold depends only on $k,\varepsilon$.
\end{corollary}

Unlike the saturated-pair bound $h>5k$, this threshold uses the additive stability constant.
It gives a stronger asymptotic density conclusion, with a much larger rank threshold.

\section{Avoiding fixed modifications in the small-cut construction}
\label{sec:high-density-protected-cuts}

The small-cut construction can avoid any fixed set of vertices.
This permits added edges to be removed after the construction.

\begin{theorem}[Protected vertices]
\label{thm:high-density-protected-cyclic-sets}
Fix $k\ge2$, $m\ge0$, and $0<\varepsilon\le1/2$.
Put $C=\lceil C_\varepsilon\rceil$ and
\[
 B_{k,m,\varepsilon}
 =(k+m+1)(9C+8k+4m+9).
\]
Let $X$ be a normalized $k$-pole of rank $h\ge B_{k,m,\varepsilon}$.
Let $Z\subseteq V(X)$ satisfy $|Z|\le m$.
If some terminal pair has at least $(1/2+\varepsilon)2^h$ paths, then $X$ contains two disjoint vertex sets $V_1,V_2$ with these properties:
\begin{enumerate}
\item Each $V_i$ avoids $Z$ and every port of $X$, and contains a cycle.
\item In every cubic completion of $X$, each $V_i$ has boundary at most four.
\end{enumerate}
In particular, either set gives a cyclic cut of size at most four in every completion.
\end{theorem}

\begin{proof}
Use the normalization and series decomposition in Theorem~\ref{thm:high-density-pair-small-cyclic-cut}.
Call the number of bridgeless pieces $q$.
Each piece retains normalized terminal-path density at least $1/2+\varepsilon$.

If $q\ge k+m+2$, exclude the first and last pieces.
Also exclude pieces containing a restored port or a vertex of $Z$ that survives normalization.
At most $(k-2)+m$ further pieces are excluded.
At least two whole pieces remain.
Use their original vertex sets, without subdivision vertices on boundary edges.
They avoid $Z$ and all ports, contain cycles, and have boundary at most two in every completion.
A protected vertex that was suppressed is one of the restored ports and was already accounted for.

If $q\le k+m+1$, some piece has rank
\[
 b\ge\frac h{k+m+1}\ge9C+8k+4m+9.
\]
Its terminal distance is at least $b-C$.
There are at most $2C$ internal layers with more than two vertices.
Mark these layers and the at most $2(k-2)$ layers meeting endpoints of subdivided edges.
Also mark the layers of protected vertices that survive normalization.
This adds at most $m$ layers.

Partition the internal distance layers into disjoint groups of four.
The number of unmarked groups is at least
\[
 \left\lfloor\frac{b-C-1}{4}\right\rfloor-2C-2(k-2)-m
 \ge6.
\]
Choose two such groups and apply Lemma~\ref{lem:four-thin-layers-short-cycle}.
The resulting cycles have length at most four, avoid $Z$, and use no subdivided edge.
Internal layers avoid the two retained terminals, and unmarked layers avoid all restored ports.
Their disjoint vertex sets remain unchanged when the ports are restored.
Each has boundary at most four in every cubic completion.
\end{proof}

\begin{corollary}[Avoiding fixed added edges]
\label{cor:high-density-forbidden-edges}
Let $X$ be a normalized $k$-pole, and let $E_0$ be at most $s$ designated edges of $X$.
If $X$ has a terminal pair of density at least $1/2+\varepsilon$ and
\[
 \beta(X)\ge(k+2s+1)(9\lceil C_\varepsilon\rceil+8k+8s+9),
\]
then the two cyclic vertex sets in Theorem~\ref{thm:high-density-protected-cyclic-sets} avoid every endpoint of every edge in $E_0$.
Deleting the edges of $E_0$ and changing incidences at their endpoints leaves both cycles and both boundaries unchanged.
The last statement applies when all other vertex incidences remain unchanged.
\end{corollary}

\begin{proof}
Protect the at most $2s$ endpoints of the designated edges.
Neither a selected vertex nor an edge incident with it is affected by the stated changes.
\end{proof}

\begin{corollary}[Unconditional augmented inequalities in a five-pole]
\label{cor:five-pole-unconditional-added-edge}
Let $Y$ be a normalized five-pole of rank $h$ inside a cubic graph $G$ with no cyclic cut of size at most four.
For distinct ports $i,j,k,l$, write
\[
 a_{kl}=2^{-h}P_{kl}(Y),\qquad
 U_{ijkl}=2^{-h}F_{\{i,j,k,l\}}(Y),\qquad
 u_{ij|kl}=2^{-h}f_{\{i,j,k,l\},ij|kl}(Y).
\]
Then, uniformly over the choices of ports,
\[
 a_{kl}+U_{ijkl}-u_{ij|kl}\le1+o(1)
 \qquad(h\to\infty).
\]
No positive lower bound on $P_{ij}(Y)$ is required.
\end{corollary}

\begin{proof}
Add an edge $ij$ to $Y$, obtaining a normalized three-pole $T$ of rank $h+1$.
The exact added-edge identity gives
\[
 P_{kl}(T)
 =2^h\bigl(a_{kl}+U_{ijkl}-u_{ij|kl}\bigr).
\]
Suppose the parenthesized expression exceeds $1+\eta$ for some fixed $\eta>0$ along arbitrarily large ranks.
Then the normalized $k$--$l$ path density in $T$ exceeds $1/2+\eta/2$.
Apply Theorem~\ref{thm:high-density-protected-cyclic-sets} with protected vertices $i,j$.
It supplies two disjoint cyclic vertex sets contained in $Y-\{i,j\}$.

The selected sets avoid $i,j$ and the three remaining ports of $T$.
Every changed incidence is therefore outside both selected sets.
Thus their boundary sizes in $G$ equal those in any cubic completion of $T$.
Both boundaries have size at most four.
Either selected set therefore gives a forbidden cyclic cut in $G$.
This contradiction proves the assertion for every fixed $\eta>0$.
There are only finitely many port choices, so the error is uniform.
\end{proof}

% Independently audited by path_stability and random_hosts.
% Exact certificates: four-regular-five-pairing-dual-checks.json.
\section{The sharp error order for cyclic cuts of size at most five}
\label{sec:five-cut-universal-crossing}

The five-port crossing bound has a short distinct-representative proof.
It uses the protected small-cut theorem and the exact five-cut identity.
No terminal-state profile classification or aggregate estimate is needed.

\begin{lemma}[Distinct representatives for five pairings]
\label{lem:five-pairing-distinct-representatives}
For each four-element subset $S\subset[5]$, choose a pairing $\alpha_S$ of $S$.
Regard $\alpha_S$ as two disjoint edges of $K_5$.
One can select an edge $e_S\in\alpha_S$ for each $S$ so that the five selected edges are distinct.
\end{lemma}

\begin{proof}
Form an auxiliary graph whose vertices are the ten edges of $K_5$.
Two auxiliary vertices are adjacent when their corresponding edges are disjoint.
This graph is the Petersen disjointness graph $KG(5,2)$.
The five chosen pairings give five distinct auxiliary edges.
Indeed, the union of the two endpoints determines the four-set $S$.

The auxiliary graph has no triangle, since three disjoint two-sets require six points.
It has no four-cycle either.
Two distinct two-sets have a union of size at least three, so there is at most one two-set disjoint from both.
Thus two distinct auxiliary vertices cannot have two common neighbors.

A five-edge subgraph of this auxiliary graph is therefore a forest, unless all five edges form one five-cycle.
For each tree component, choose a root and assign every edge to its endpoint farther from the root.
For a five-cycle, orient the cycle and assign each edge to its head.
These assignments use distinct auxiliary vertices.
They give the required distinct representative edges of $K_5$.
\end{proof}

\begin{lemma}[A universal five-port functional]
\label{lem:universal-five-port-functional}
Let $a_e\ge0$ be indexed by the ten pairs of $[5]$.
For each four-set $S$, let $u_{S,\alpha}\ge0$ be indexed by its three pairings, and put $U_S=\sum_\alpha u_{S,\alpha}$.
Suppose
\[
 a_e\le\tfrac12+\eta_1,
 \qquad
 a_e+U_S-u_{S,\alpha}\le1+\eta_2
 \quad(e\in\alpha),
\]
where $\eta_1,\eta_2\ge0$.
Then
\[
 \sum_e a_e+\sum_S\max_\alpha(U_S-u_{S,\alpha})
 \le\frac{15}{2}+5\eta_1+5\eta_2.
\]
\end{lemma}

\begin{proof}
For each $S$, choose a maximizing pairing $\alpha_S$.
Apply Lemma~\ref{lem:five-pairing-distinct-representatives} and choose distinct $e_S\in\alpha_S$.
Summing the five corresponding augmented inequalities gives
\[
 \sum_S\bigl(U_S-u_{S,\alpha_S}\bigr)
 +\sum_S a_{e_S}\le5+5\eta_2.
\]
Exactly five path coordinates remain.
Their total is at most $5/2+5\eta_1$.
Adding them proves the stated inequality.
\end{proof}

\begin{proposition}[A uniform deficit across a five-edge cut]
\label{prop:five-cut-universal-deficit}
Let $G$ have no cyclic edge cut of size at most four.
Suppose $G$ is obtained by joining the ports of normalized five-poles $X,Y$ of ranks $h,k$.
If $k\to\infty$, then their normalized crossing contribution is at most
\[
 \frac{15}{2}+o(1).
\]
The estimate is uniform in $X$ and its rank $h$.
\end{proposition}

\begin{proof}
Normalize the states of $Y$ by $2^k$ and write them as $a_e,u_{S,\alpha},U_S$.
Corollaries~\ref{cor:cyclic-five-path-density} and~\ref{cor:five-pole-unconditional-added-edge} give the hypotheses of Lemma~\ref{lem:universal-five-port-functional}, with errors tending to zero as $k\to\infty$.

Normalize the path coordinates of $X$ by $2^h$ and write them as $b_e$.
They all lie in $[0,1]$.
Write its normalized pairing coordinates as $v_{S,\alpha}\ge0$.
For each $S$, their sum is at most one.
The exact crossing contribution is
\[
 \mathcal B(X,Y)
 =\sum_e b_ea_e
  +\sum_S\sum_\alpha v_{S,\alpha}
       (U_S-u_{S,\alpha}).
\]
Every coefficient $U_S-u_{S,\alpha}$ is nonnegative.
The coordinate bounds on $X$ therefore give
\[
 \mathcal B(X,Y)
 \le\sum_ea_e+\sum_S\max_\alpha(U_S-u_{S,\alpha})
 \le\frac{15}{2}+o(1).
\]
All errors concern only the rank of $Y$.
This proves the stated uniformity.
\end{proof}

\begin{theorem}[Sharp-order five-cut error]
\label{thm:sharp-order-five-cut-error}
There is an absolute constant $K_5$ such that every connected loopless cubic multigraph of rank $r$ with a cyclic edge cut of size at most five satisfies
\[
 c(G)\le2^{r-1}+K_5r2^{r/2}.
\]
The factor $r$ cannot be replaced by $o(r)$ throughout this class.
Moreover, for all sufficiently large $r$, a graph with a cyclic five-edge cut and no smaller cyclic cut satisfies
\[
 c(G)<2^{r-1}.
\]
\end{theorem}

\begin{proof}
If $G$ has a cyclic cut of size at most four, apply Theorem~\ref{thm:sharp-order-four-cut-error}.
Suppose no such cut exists, and choose a cyclic five-edge cut.

Both shores are connected.
Indeed, any component containing a cycle must have at least five boundary edges, since its complement contains the opposite shore's cycle.
It uses all five cut edges.
Any other component on its shore would have no edge to its complement, contrary to the connectedness of $G$.

No vertex has two or more incidences in the cut.
Such a vertex has internal degree at most one.
Moving it to the opposite shore preserves a cycle on both sides and reduces the boundary to at most four.
This would contradict the hypothesis.
Thus the two shores are normalized five-poles $X,Y$.
Write their ranks as $h,k\ge1$ and assume $h\le k$.
Then
\[
 r=h+k+4,\qquad k\ge(r-4)/2.
\]

The uniform crossing bound gives
\[
 \frac{c(G)}{2^{h+k}}
 \le\frac{15}{2}+o(1)
    +\frac{c(X)+c(Y)}{2^{h+k}}.
\]
The affine-space bound gives $c(X)\le2^h-1$.
The asymptotic cycle bound gives $c(Y)\le(1/2+o(1))2^k$.
Since $h\ge1$, their normalized total is at most
\[
 2^{-k}+(1/2+o(1))2^{-h}
 \le\frac14+o(1).
\]
Consequently,
\[
 \frac{c(G)}{2^{h+k}}\le\frac{31}{4}+o(1)<8
\]
for all sufficiently large $r$.
This proves the strict one-half bound in the absence of smaller cyclic cuts.

The remaining bounded ranks are absorbed into $K_5$, using $c(G)\le2^r-1$.
The cyclic-square examples already have a cyclic four-edge cut and excess $\Theta(r2^{r/2})$.
They also belong to the present class, so the polynomial factor is sharp.
\end{proof}

The unresolved sharp-order case is now confined to cubic kernels with no cyclic edge cut of size at most five.
The unrestricted polynomial-error theorem applies to those kernels, but does not give the half-exponent error there.


\section{Verification record}
The exact four-cut checker tests 600 gluings, 120 added-edge and capping identities, and 768 individual forest mappings.
It compares the state formula with independent global cycle-space enumeration.
A separate checker verifies the cyclic-square total cycle formula for two through seven blocks and the disjoint-cap packing formula for one through six square blocks.
The layer encodings and terminal-path saturation statements have additional independent checks recorded in the shared research directory.
These finite checks support the identities; the asymptotic bounds follow from the proofs.

The five-cut gluing identity has 1,920 independent exact checks.
All 243 choices in the Petersen-graph representative argument also pass exact integer checks.
The local four-layer lemma has an exhaustive check of its finite strip configurations.

The unrestricted $O(r2^{r/2})$ question is now reduced, within the established cubic-kernel reductions, to cyclically six-edge-connected kernels.
The six-terminal state analysis remains work in progress.
\begin{thebibliography}{9}
\bibitem{HK} P. Horak and D. Kim, \emph{Connected cubic graphs with the maximum number of perfect matchings}, Journal of Graph Theory 99 (2022), 671--690. Updated author version: \url{https://arxiv.org/abs/2006.13459}.
\end{thebibliography}
\end{document}
